% Generated by tools/make_body.py; edit generator, not this file.
% UNPROOFREAD OCR TRANSCRIPTION: page-image verification remains required.
\chapter{Global Analytic Functions}\label{chap:8}

% Original PDF physical page 299
\sourcepage{299}
\setcounter{section}{0}\section{Analytic Continuation}\index{Function!analytic!global|(}\index{Global analytic function|(}\label{sec:8.1}
In the preceding chapters we have stressed that all functions must be well defined and, therefore, single-valued. In the case of functions such as $\sqrt{z}$ and $\log z$ which are not uniquely determined by their analytic expressions, a special effort was needed to show that, under favorable circumstances, a single-valued branch can be selected. While this answers the need for logical clarity, it does not do justice to the fact that the ambiguity of the square root or the logarithm is an essential feature which cannot be ignored. There is thus a clear need for a concept that emphasizes rather than circumvents multiple valuedness.

\setcounter{section}{1}\setcounter{subsection}{0}\subsection{The Weierstrass Theory}\label{sec:8.1.1}
Weierstrass\index{Weierstrass K@Weierstrass, K.|(}\index{Weierstrass K@Weierstrass, K.}, in contrast to Riemann, who favored a more geometric outlook, wanted to build the whole theory of analytic functions from the concept of power series. For Weierstrass the basic building block was a power series
\begin{equation*}
P (z - \zeta) = a _ {0} + a _ {1} (z - \zeta) + \dots + a _ {n} (z - \zeta) ^ {n} + \dots
\end{equation*}
with a positive radius of convergence $r(P)$. Such a series is determined by a complex number $\zeta$, the center of the power series, and a sequence $\{a_n\}_0^\infty$ of complex coefficients. The radius of convergence is given by Hadamard's formula $r(P)^{-1} = \varlimsup_{n\to \infty}|a_n|^{1 / n}$. It is an essential requirement that $r(P) > 0$, for only then does the power series define an analytic function $f(z)$ in the disk $D = \{z||z - \zeta| < r(P)\}$.


% Original PDF physical page 300
\sourcepage{300}
Given a point $\zeta_1 \in D$, the function $f(z)$ has a Taylor development $P_1(z - \zeta_1)$ about $\zeta_1$. It converges in a disk $D_1$ whose radius $r(P_1)$ is at least equal to $r(P_0) - |\zeta_1 - \zeta|$, but may be larger. The new series defines an analytic function $f_1(z)$ in $D_1$ which is said to be obtained from $f(z)$ by direct analytic continuation\index{Analytic continuation!direct}\index{Analytic continuation}. Together, $f$ and $f_1$ define an analytic function in $D \cup D_1$, for they are equal in the intersection $D \cap D_1$. If $D_1$ is not contained in $D$, the new function is an extension of $f$ to a larger region, and that is the purpose of the construction.

This process can be repeated any number of times. In the general case we have to consider a succession of power series $P_{0}(z - \zeta_{0})$, $P_{1}(z - \zeta_{1})$, $\ldots$, $P_{n}(z - \zeta_{n})$, each of which is a direct analytic continuation of the preceding one. In other words, if $P_{k}$ converges to a function $f_{k}$ in the disk $D_{k}$, then $\zeta_{k} \in D_{k-1}$ and $f_{k} = f_{k-1}$ in $D_{k-1} \cap D_{k}$. It does not follow that $f_{0}, \ldots, f_{n}$ define a single-valued function in $D_{0} \cup D_{1} \cup \cdots \cup D_{n}$, for if $D_{k}$ meets a $D_{h}$ with h different from k - 1 and $k + 1$, there is no guarantee that $f_{h} = f_{k}$ in $D_{h} \cap D_{k}$. It is precisely this possibility that leads beyond the notion of function in the strict sense of having only one value at each point of its domain.

As soon as there exist power series $P_{0}, P_{1}, \ldots, P_{n}$ as above, one says that $P_{n}$ is an analytic continuation of $P_{0}$. Weierstrass\index{Weierstrass K@Weierstrass, K.} considers the totality of all power series $P(z - \zeta)$ that can be obtained from $P_{0}(z - \zeta_{0})$ by analytic continuation. This set of power series will be called an analytic function in the sense of Weierstrass.

The property of one power series to be an analytic continuation of another is evidently an equivalence relation. An analytic function in the sense of Weierstrass is nothing but an equivalence class with respect to this relation, and the initial power series $P_{0}$ is in no distinguished position within its class. The underlying idea is that two power series which belong to the same equivalence class are different forms of the same function.

\setcounter{section}{1}\setcounter{subsection}{1}\subsection{Germs and Sheaves}\index{Germ}\label{sec:8.1.2}
The Weierstrass theory has mostly historical interest, for the restriction to power series and their domains of convergence is more of a hindrance than a help. It should, nevertheless, be recognized that the idea of Weierstrass is still the basis for our understanding of multiple-valuedness in the theory of complex analytic functions.

We shall outline a more direct approach which is more in line with the somewhat sophisticated ideas that dominate the recent theory of analytic functions of several complex variables. Because of the limited scope of this book we have to be content to borrow some of the terminology and use it to simplify some proofs.

An analytic function $f$ defined in a region $\Omega$ will constitute a function element\index{Function element}, denoted by $(f, \Omega)$, and a global analytic\index{Function!analytic!global} function\index{Global analytic function} will appear as a
% Original PDF physical page 301
\index{Weierstrass K@Weierstrass, K.|)}\sourcepage{301}
collection\index{Projection of germ|(} of function elements which are related to each other in a prescribed manner.

Two function elements $(f_{1},\Omega_{1})$ and $(f_{2},\Omega_{2})$ are said to be direct analytic continuations of each other if $\Omega_{1}\cap\Omega_{2}$ is nonempty and $f_{1}(z)=f_{2}(z)$ in $\Omega_{1}\cap\Omega_{2}$. More specifically, $(f_{2},\Omega_{2})$ is called a direct analytic continuation of $(f_{1},\Omega_{1})$ to the region $\Omega_{2}$. There need not exist any direct analytic continuation to $\Omega_{2}$, but if there is one, it is uniquely determined. For suppose that $(f_{2},\Omega_{2})$ and $(g_{2},\Omega_{2})$ are two direct analytic continuations of $(f_{1},\Omega_{1})$; then $f_{2}=g_{2}$ in $\Omega_{1}\cap\Omega_{2}$, and because $\Omega_{2}$ is connected, this implies $f_{2}=g_{2}$ throughout $\Omega_{2}$. We note that if $\Omega_{2}\subset\Omega_{1}$, then the direct analytic continuation of $(f_{1},\Omega_{1})$ is $(f_{1},\Omega_{2})$.

As in the case of power series we consider chains $(f_{1},\Omega_{1})$, $(f_{2},\Omega_{2})$, $\ldots$, $(f_{n},\Omega_{n})$ such that $(f_{k},\Omega_{k})$ and $(f_{k+1},\Omega_{k+1})$ are direct analytic continuations of each other, and we say that $(f_{n},\Omega_{n})$ is an analytic continuation of $(f_{1},\Omega_{1})$. This defines an equivalence relation, and the equivalence classes are called global analytic\index{Function!analytic!global} functions\index{Global analytic function}. As a typographical device the global analytic function determined by the function element $(f,\Omega)$ will be denoted by bold type, $\mathbf{f}$. For a more flexible terminology $(f,\Omega)$ is also referred to as a branch\index{Branch} of $\mathbf{f}$. While $(f,\Omega)$ determines $\mathbf{f}$ uniquely, the converse is not true; $\mathbf{f}$ may have several branches over the same $\Omega$.

It is quite obvious that global analytic functions can be identified with analytic functions in the sense of Weierstrass, and we have gained almost nothing in generality. There is, however, a more fruitful point of view. Instead of pairs $(f,\Omega)$ we shall consider pairs $(f,\zeta)$ where $\zeta$ is a point and f is analytic at $\zeta$, that is to say, f is defined and analytic in some open set that contains $\zeta$. Two pairs $(f_{1},\zeta_{1})$ and $(f_{2},\zeta_{2})$ shall be equivalent if and only if $\zeta_{1} = \zeta_{2}$ and $f_{1} = f_{2}$ in some neighborhood of $\zeta_{1}$. The conditions for an equivalence relation are obviously fulfilled. The equivalence classes are called germs\index{Function!analytic!germ of}, or more specifically germs of analytic functions. Each germ determines a unique $\zeta$, the projection of the germ, and we use the notation $\mathbf{f}_{\zeta}$ to indicate a germ with projection $\zeta$. A function element $(f,\Omega)$ gives rise to a germ $\mathbf{f}_{\zeta}$ for each $\zeta \in \Omega$; conversely, every $\mathbf{f}_{\zeta}$ is determined by some $(f,\Omega)$.

The reader will of course recognize that the germs $\mathbf{f}_{\zeta}$ can be identified with the corresponding convergent power series $P(z - \zeta)$, and we are back where we started. However, by introducing the notion of germ we have isolated an essential property of convergent power series, namely the fact that two power series with the same center are identical if and only if they represent the same function in some neighborhood of the center. In pursuit of this idea it becomes clear that we could equally well consider germs of other classes of functions, for instance, germs of continuous functions, germs of functions of class $C^{k}$, etc., for which the identification with power series is no longer possible. Although we are mainly interested in
% Original PDF physical page 302
\index{Projection of germ|)}\sourcepage{302}
germs of analytic functions, we are nevertheless going to take a slightly more general point of view.

Let $D$ be an open set in the complex plane. The set of all germs $\mathbf{f}_{\zeta}$ with $\zeta \in D$ is called a sheaf\index{Sheaf} over $D$; we shall denote it by $\mathfrak{S}$ or $\mathfrak{S}_D$. If we are dealing with germs of analytic functions, $\mathfrak{S}_D$ is called the sheaf of germs of analytic functions over $D$. There is a projection map $\pi: \mathfrak{S} \to D$ which maps $\mathbf{f}_{\zeta}$ on $\zeta$. For a fixed $\zeta \in D$ the inverse image $\pi^{-1}(\zeta)$ is called the stalk\index{Stalk} over $\zeta$; it is denoted by $\mathfrak{S}_{\zeta}$.

The set $\mathfrak{S}$ is interesting because it carries a twofold structure: one topological and one algebraic. First, $\mathfrak{S}$ can be made into a topological space, which enables us to speak of continuous mappings. Second, there is an obvious algebraic structure on each stalk, for it is clear what we mean by $\mathbf{f}_{\zeta} + \mathbf{g}_{\zeta}$ or $\mathbf{f}_{\zeta} \cdot \mathbf{g}_{\zeta}$. For the sake of simplicity we shall fix our attention on the additive structure. In terms of this structure each stalk is an abelian group.

We are ready for a fairly general definition.

\begin{originalstatement}{Definition}{1}\label{def:8.1}
A sheaf over $D$ is a topological space $\mathfrak{S}$ and a mapping $\pi\colon\mathfrak{S}\to D$ with the following properties:
(i) The mapping $\pi$ is a local homeomorphism; this shall mean that each $s \in \mathfrak{S}$ has an open neighborhood $\Delta$ such that $\pi(\Delta)$ is open and the restriction of $\pi$ to $\Delta$ is a homeomorphism.

(ii) For each $\zeta \in D$ the stalk $\pi^{-1}(\zeta) = \mathfrak{S}_{\zeta}$ has the structure of an abelian group.

(iii) The group operations are continuous in the topology of $\mathfrak{S}$.
\end{originalstatement}

Actually, D can be an arbitrary topological space, but we shall think of D as an open set in the complex plane. Also, the structure of an abelian group can be replaced by other algebraic structures.

We shall now verify that the sheaf S of germs of analytic functions satisfies the conditions in Definition~\hyperref[def:8.1]{1}. For this purpose we must first introduce a topology on S. It is awkward, and unnecessary, to make S a metric space. Instead, we need merely specify the subsets of S which are to be the open sets in the topology. Our characterization of open sets shall be as follows: A set $V \subset S$ is open if for every $s_{0} \in V$ there exists a function element $(f,\Omega)$ such that (1) $\pi(s_{0}) = \zeta_{0} \in \Omega$, (2) $(f,\Omega)$ determines the germ $s_{0}$ at $\zeta_{0}$, (3) all the germs $\mathbf{f}_{\zeta}$ determined by $(f,\Omega)$ are in $V$. The reader will have no difficulty verifying that the conditions of Chap.~\hyperref[chap:3]{3}, Def. 8, are satisfied.

With $s_0$ and $(f, \Omega)$ as above, let $\Delta$ be the set of all the germs $\mathbf{f}_{\zeta}$ determined by $(f, \Omega)$. Owing to our definition of open set, it is quite obvious that $\Delta$ is an open neighborhood of $s_0$, and that the mapping $\pi: \Delta \to \Omega$ is a homeomorphism. Thus condition (i) of the definition is fulfilled.


% Original PDF physical page 303
\sourcepage{303}
Condition (ii) needs no proof. Condition (iii) is also easy, but it is important to understand what is involved. Addition and subtraction make sense only for germs on the same stalk; it is sufficient to consider subtraction. Consider two germs $s_0, s_0'$ with $\pi(s_0) = \pi(s_0') = \zeta_0$. Let them be determined by function elements $(f, \Omega)$ and $(g, \Omega)$ with $\zeta_0 \in \Omega$; for the sake of simplicity we have chosen the same $\Omega$ for both function elements. If $s \in \Delta_0, s' \in \Delta_0'$ with $\pi(s) = \pi(s') = \zeta$, then $s - s'$ is the germ determined by $(f - g, \Omega)$ at $\zeta$. When $\zeta$ ranges over $\Omega, s - s'$ ranges over a neighborhood of $s_0 - s_0'$; moreover, $\pi(s - s') = \pi(s) - \pi(s')$. The projection maps establish homeomorphisms between $\Delta, \Delta_0, \Delta_0'$, and $\Omega$. It is therefore clear that we can shrink $\Delta_0$ and $\Delta_0'$ so as to make $\Delta$ contained in any prescribed neighborhood of $s_0 - s_0'$, thereby proving the continuity.

\setcounter{section}{1}\setcounter{subsection}{2}\subsection{Sections and Riemann Surfaces}\index{Riemann surface|(}\index{Riemann surface}\label{sec:8.1.3}
Let $\mathfrak{S}$ be a sheaf over $D$ and consider an open set $U \subset D$. A continuous mapping $\varphi: U \to \mathfrak{S}$ is called a section\index{Section|(}\index{Section} over $U$ if the composed mapping $\pi \circ \varphi$ is the identity mapping of $U$ on itself. It follows from this condition that $\varphi(\zeta_1) = \varphi(\zeta_2)$ implies $\zeta_1 = \zeta_2$; hence $\varphi$ is one to one, and its inverse is $\pi$ restricted to $\varphi(U)$. Thus every section is a homeomorphism.

Every point $s_0 \in \mathfrak{S}$ is in the image $\varphi(U_0)$ of some section; we need only take $U_0 = \pi(\Delta)$ where $\Delta$ is the neighborhood whose existence is postulated in (ii), and $\varphi$ equal to the inverse of $\pi$ as restricted to $\Delta$.

The set of all sections over a fixed U is denoted by $\Gamma(U,\mathfrak{S})$. If nonempty, it has the structure of an abelian group, for it makes sense to define $\varphi - \psi$ as the section with values $\varphi(\zeta) - \psi(\zeta)$. Let $0_{\zeta}$ be the zero element of the stalk $S_{\zeta}$, and define a function $\omega$ by setting $\omega(\zeta) = 0_{\zeta}$. We claim that $\omega$ is continuous, and hence a section; it is called the zero section, and it acts as a zero element for the group $\Gamma(U,\mathfrak{S})$.

To prove the continuity, consider a point $\zeta_0 \in U$ and an $s_0 \in \mathfrak{S}_{\zeta_0}$ (for instance, $0_{\zeta_0}$). According to our previous remark $s_0$ is in some $\varphi(U_0)$. By condition (iii) $\varphi - \varphi = \omega$ is continuous in $U_0$. Since $\zeta_0$ is arbitrary, $\omega$ is continuous on all of $U$, and hence a section. We have shown that the zero section always exists, and $\Gamma(U, \mathfrak{S})$ is not empty. From now on the zero section will be denoted by 0.

If $U$ is connected and $\varphi, \psi \in \Gamma(U, \mathfrak{S})$, then either $\varphi$ and $\psi$ are identical, or the images $\varphi(U)$ and $\psi(U)$ are disjoint. Indeed, the sets with $\varphi - \psi = 0$ and $\varphi - \psi \neq 0$ are both open.

We have carried out this discussion in some detail to show how the postulates work. The special case of the sheaf of germs of analytic functions is rather trivial, for in that case $\Gamma(U,\mathfrak{S})$ can be interpreted as the additive group of analytic (“single-valued”) functions on $U$. The zero section is nothing but the constant 0.

In what follows $\mathfrak{S}$ will always be the sheaf of germs of analytic func%
% Original PDF physical page 304
\sourcepage{304}%
tions\index{Functional relation|(} over the whole complex plane. The components of $\mathfrak{S}$, regarded as a topological space, can be identified with the global analytic\index{Function!analytic!global} functions\index{Global analytic function}. To see this, let $s_0 \in \mathfrak{S}$ be a germ determined by the function element $(f_0, \Omega_0)$, and let $(f_1, \Omega_1)$ be a direct analytic continuation of $(f_0, \Omega_0)$; we remind the reader that $\Omega_0$ and $\Omega_1$ are assumed to be connected. Because $f_0 = f_1$ in $\Omega_0 \cap \Omega_1$ the sets $\Delta_0$ and $\Delta_1$ of germs determined by these two function elements intersect; as homeomorphic images of $\Omega_0, \Omega_1$ the sets $\Delta_0, \Delta_1$ are connected, and the same is consequently true of their union, $\Delta_0 \cup \Delta_1$. It follows that all the function elements that are obtainable from $(f_0, \Omega_0)$ by a chain of direct analytic continuations give rise to germs contained in the component $\mathfrak{S}_0$ of $s_0$. On the other hand, let $\mathfrak{S}_0'$ be the set of germs in $\mathfrak{S}_0$ which can be determined by an analytic continuation $(f, \Omega)$ of $(f_0, \Omega_0)$. It is readily seen that $\mathfrak{S}_0'$ and its complement in $\mathfrak{S}_0$ are both open. Hence $\mathfrak{S}_0' = \mathfrak{S}_0$, and we conclude that $\mathfrak{S}_0$ consists precisely of all the germs belonging to a global analytic function.

In spite of this identification, it is more suggestive to regard $\mathfrak{S}_{0}$ as the domain of the global analytic function, which we shall now denote by $\mathbf{f}$, its value at $\mathbf{f}_{\zeta}$ being nothing but the constant term in the power series associated with that germ. With this interpretation $\mathfrak{S}_{0}$ is referred to as the Riemann surface\index{Riemann surface} of $\mathbf{f}$. It is indeed quite similar to the elementary Riemann surfaces which were briefly studied in Chap.~\hyperref[chap:3]{3}, Sec.~\hyperref[sec:3.4.3]{4.3}, and it serves the same purpose, namely to make $\mathbf{f}$ single-valued. One can picture $\mathfrak{S}_{0}$ as being spread out in layers over the plane, and the sheets, if that is what one wants to call them, are images of sections. It should be noticed that we are not yet including the branch points, whose role will be investigated later.

For greater clarity, let the Riemann surface of a global analytic function $\mathbf{f}$ be denoted by $\mathfrak{S}_0(\mathbf{f})$. Given two global functions $\mathbf{f}$ and $\mathbf{g}$, there may exist a mapping $\theta: \mathfrak{S}_0(\mathbf{f}) = \mathfrak{S}_0(\mathbf{g})$ such that (1) $\pi \circ \theta = \pi$, and (2) $\theta$ is a local homeomorphism. In these circumstances $\mathbf{g} \circ \theta$ is a single-valued function on $\mathfrak{S}_0(f)$; usually, the notation is simplified and one agrees to write $\mathbf{g}$ instead of $\mathbf{g} \circ \theta$. In this way all the derivatives $\mathbf{f}', \mathbf{f}'', \ldots$ are defined on the Riemann surface of $\mathbf{f}$. All entire functions $\mathbf{h}$ are automatically defined on every $\mathfrak{S}_0(\mathbf{f})$, and if $\mathbf{g}, \mathbf{h}, \ldots$ are defined on $\mathfrak{S}_0(\mathbf{f})$, so is every polynomial $G(\mathbf{f}, \mathbf{g}, \mathbf{h}, \ldots)$.

There is a classical principle known as the permanence of functional relations. Suppose that certain function elements $(f,\Omega)$, $(g,\Omega)$, $(h,\Omega)$, $\ldots$ can be continued analytically whenever $(f,\Omega)$ can be continued, directly or through a chain of direct continuations. Assume moreover that $G(f,g,h,\ldots)=0$ on $\Omega$. Then the same relation holds for all analytic continuations, a fact that may be expressed by $G(\mathbf{f},\mathbf{g},\mathbf{h},\ldots)=0$. In particular, if a germ satisfies a polynomial differential equation $G(z,f,f',\ldots,f^{(n)})=0$, then the global function $\mathbf{f}$ satisfies the same equation.


% Original PDF physical page 305
\index{Functional relation|)}\index{Riemann surface|)}\index{Section|)}\sourcepage{305}
\setcounter{section}{1}\setcounter{subsection}{3}\subsection{Analytic Continuation along Arcs}\label{sec:8.1.4}
Let $\gamma: [a,b] \to \mathbf{C}$ be an arc in the complex plane. Consider a global analytic\index{Function!analytic!global} function\index{Global analytic function} $\mathbf{f}$ and its Riemann surface $\mathfrak{S}_0(\mathbf{f})$, defined, as before, to be a component of the sheaf $\mathfrak{S}$ of all germs of analytic functions. An arc $\bar{\gamma}: [a,b] \to \mathfrak{S}_0(\mathbf{f})$ is said to be an analytic continuation of $\mathbf{f}$ along $\gamma$ if $\pi \circ \bar{\gamma} = \gamma$, i.e., if $\bar{\gamma}(t)$ projects on $\gamma(t)$ for all $t \in [a,b]$. Naturally, by the definition of arc, $\bar{\gamma}(t)$ must be continuous on $[a,b]$ in the topology of $\mathfrak{S}_0(\mathbf{f})$. In another terminology, $\bar{\gamma}$ is also called a lifting of $\gamma$ to $\mathfrak{S}_0(\mathbf{f})$.

Continuation along an arc corresponds to the intuitive notion of a continuously changing germ. The existence of a continuation is not guaranteed, but the following important uniqueness theorem is valid:

\begin{originalstatement}{Theorem}{1}\label{thm:8.1}
Two analytic continuations $\bar{\gamma}_1$ and $\bar{\gamma}_2$ of a global analytic function $\mathbf{f}$ along the same arc $\gamma$ are either identical, or $\bar{\gamma}_1(t) \neq \bar{\gamma}_2(t)$ for all $t$.
\end{originalstatement}
The proof is a triviality. Because $\pi$ is a local homeomorphism the image of $\bar{\gamma}_{1}-\bar{\gamma}_{2}$ cannot contain a point of the zero section without being contained in it.

By virtue of this theorem a continuation is uniquely determined by its initial value, the germ $\bar{\gamma}(a)$; the initial germ is of the form $\mathbf{f}_{\gamma(a)}$, but $\mathbf{f}$ may have several germs of this form. Once the initial germ is specified we have the right to speak of the analytic continuation from that germ, provided that such a continuation exists.

It may well happen that $\mathbf{f}$ does not have a continuation along $\gamma$, or that a continuation exists for some initial germs, but not for all. Let us investigate the case of an initial germ $\mathbf{f}_{\gamma(a)}$ which cannot be continued along $\gamma$. If $t_0 > a$ is sufficiently close to $a$, there will always exist a continuation of the initial germ along the subarc of $\gamma$ that corresponds to the interval $[a, t_0]$; indeed, if $\mathbf{f}_{\gamma(a)}$ is determined by the function element $(f_0, \Omega_0)$, this is trivially the case if the subarc is contained in $\Omega_0$. The least upper bound of all such $t_0$ is a number $\tau$ with $a < \tau < b$, and the continuation will be possible for $t_0 < \tau$, impossible for $t_0 \geq \tau$. In a certain sense the subarc $\gamma[a, \tau]$ leads to a point at which $\mathbf{f}$ ceases to be defined. This subarc is called a singular path\index{Singular path} from the given initial germ; less precisely, it is said to lead to a singular point\index{Singular point} over $\gamma(\tau)$. Observe that when $t$ approaches $\tau$ from below, the radius of convergence of the power series representing the germ $\bar{\gamma}(t)$ will tend to zero.

The connection between continuation along arcs and stepwise continuation by means of a chain of direct analytic continuations requires further illumination. In the first place, if $(f_{1},\Omega_{1})$, $(f_{2},\Omega_{2})$, $\ldots$, $(f_{n},\Omega_{n})$ is a chain of direct analytic continuations, it is always possible to connect a point $\zeta_{1}\in\Omega_{1}$ to a point $\zeta_{n}\in\Omega_{n}$ by means of an arc $\gamma$ such that there exists a continuation $\bar{\gamma}$ with initial germ $(f_{1},\zeta_{1})$ and terminal germ $(f_{n},\zeta_{n})$. Indeed,
% Original PDF physical page 306
\sourcepage{306}
it is sufficient to let $\gamma$ be composed of a subarc $\gamma_{1}$ in $\Omega_{1}$ from $\zeta_{1}$ to a point $\zeta_{2} \in \Omega_{1} \cap \Omega_{2}$, a second subarc $\gamma_{2}$ in $\Omega_{2}$ from $\zeta_{2}$ to $\zeta_{3} \in \Omega_{2} \cap \Omega_{3}$, and so on. The continuation along $\gamma$ is defined by $\bar{\gamma}(t) = (f_{k}, \zeta(t))$ on $\gamma_{k}$.

Conversely, if $\bar{\gamma}$ is given, we can find a chain of direct analytic continuations which follows the arc $\gamma$ in the same way as in the preceding construction. In fact, by Heine-Borel's lemma the parametric interval $[a,b]$ can be subdivided into $[a,t_1],[t_1,t_2],\ldots,[t_{n-1},b]$ such that $\bar{\gamma}(t) = (f_k,\gamma(t))$ in $[t_{k-1},t_k]$ for suitably chosen function elements $(f_k,\Omega_k)$. Although $(f_k,\Omega_k)$ and $(f_{k+1},\Omega_{k+1})$ need not be direct analytic continuations of each other, they are at least direct continuations of their common restrictions to a neighborhood of $\gamma(t_k)$.

In order to illustrate the use of continuations along arcs we shall define the logarithm as a global analytic\index{Function!analytic!global} function\index{Global analytic function}. For this purpose we want to show that the set of all function elements $(f,\Omega)$ with $e^{f(\zeta)} = \zeta$ in $\Omega$ is a global analytic function.

We need only make sure that any two function elements $(f_{1},\Omega_{1})$, $(f_{2},\Omega_{2})$ in this collection can be joined by a chain of direct analytic continuations, for the permanence of functional relations will guarantee that the intermediate function elements belong to the same collection. Choose points $\zeta_{1}\in\Omega_{1}$, $\zeta_{2}\in\Omega_{2}$ and join them by an arc $\gamma(t)$, $t\in[a,b]$ which does not pass through the origin; this is possible because neither $\zeta_{1}$ nor $\zeta_{2}$ can be zero. Consider the function
\begin{equation*}
\varphi (t) = f _ {1} (\zeta_ {1}) + \int_ {a} ^ {t} \frac {\gamma^ {\prime} (t)}{\gamma (t)} d t.
\end{equation*}
By differentiation, $\gamma(t)e^{-\varphi(t)}$ is a constant; for $t = a$ the value is 1, and hence $e^{\varphi(t)} = \gamma(t)$. For a given $t$ there exists, for instance in the disk $\Omega = \{\zeta | |\zeta - \gamma(t)| < |\gamma(t)|\}$, a uniquely determined branch $f(\zeta)$ of $\log \zeta$ which takes the value $\varphi(t)$ for $\zeta = \gamma(t)$. It is clear that $\bar{\gamma}(t)$ will be a continuation along $\gamma$. The germ $\bar{\gamma}(b)$ at the end point may not coincide with the one determined by $(f_2, \Omega_2)$, but its value at $\zeta_2$ will differ from $f_2(\zeta_2)$ by a multiple of $2\pi i$. In order to obtain the right value at $\zeta_2$, all that remains is to continue from $\bar{\gamma}(b)$ along a closed curve which circles the origin the right number of times. Finally, the arcwise continuation can be replaced by a finite chain of direct analytic continuations, and we have shown that our construction defines the logarithm as a global analytic function.

\subsection*{EXERCISES}
\originalnumber{1}{exercise:8.1.4.1.1}
1. If a function element is defined by a power series inside its circle of convergence, supposed to be of finite radius, prove that at least one radius is a singular path for the global analytic function which it determines. ("A power series has at least one singular point on its circle of convergence.")


% Original PDF physical page 307
\sourcepage{307}
\originalnumber{2}{exercise:8.1.4.1.2}
2.\index{Homotopy|(} If a function element $(f,\Omega)$ has no direct analytic continuations other than the ones obtained by restricting $f$ to a smaller region, then the boundary of $\Omega$ is called a natural boundary\index{Natural boundary} for $f$. Prove that the series $\sum_{n=0}^{\infty}z^{n!}$ has the unit circle as a natural boundary. Hint: Show that the function tends to infinity on every radius whose argument is a rational multiple of $\pi$.

\originalnumber{3}{exercise:8.1.4.1.3}
3. Show that the function $\lambda(\tau)$ introduced in Chap.~\hyperref[chap:7]{7}, Sec.~\hyperref[sec:7.3.4]{3.4}, has the real axis as a natural boundary.

\setcounter{section}{1}\setcounter{subsection}{4}\subsection{Homotopic Curves}\label{sec:8.1.5}
We must now study the topological properties of closed curves in a region from a point of view which is fundamental for the theory of analytic continuations. The question which interests us is the behavior of an arc under continuous deformations. From an intuitive standpoint this is an extremely simple notion. If $\gamma_{1}$ and $\gamma_{2}$ are two arcs with common end points, contained in a region $\Omega$, it is very natural to ask whether $\gamma_{1}$ can be continuously deformed into $\gamma_{2}$ when the end points are kept fixed and the moving arc is confined to $\Omega$. For instance, in Fig.~\hyperref[fig:8-1]{8-1} the arc $\gamma_{1}$ can be deformed into $\gamma_{2}$, but not into $\gamma_{3}$. Two arcs which can be deformed into each other are said to be homotopic in $\Omega$. This is evidently an equivalence relation.

A precise definition must of course be given. Fortunately, the physical conception of deformation\index{Deformation} has an almost immediate interpreta%
\bookfigure{307}{1}{FIG. 8-1. Homotopic arcs.}%
% Original PDF physical page 308
\sourcepage{308}%
tion in mathematical terms. It is indeed clear that a deformation of an arc can be described by means of a continuous function $\gamma(t,u)$ of two variables, the point $(t,u)$ ranging over a rectangle $[a,b] \times [0,1]$ (Fig.~\hyperref[fig:8-2]{8-2}). To every fixed value $u = u_0$ there corresponds an arc $\gamma(t,u_0)$, and the effect of the deformation is to change the initial arc $\gamma(t,0)$ into $\gamma(t,1)$. The deformation takes place within $\Omega$ if $\gamma(t,u) \in \Omega$ for all $(t,u)$, and it is a deformation with fixed end points if $\gamma(a,u)$ and $\gamma(b,u)$ are constant. To every fixed value $t = t_0$ there corresponds an arc $\gamma(t_0,u)$, $u \in [0,1]$, which may be called a deformation path.

We are led to the following formal definition of homotopy\index{Homotopy}:

\begin{originalstatement}{Definition}{2}\label{def:8.2}
Two arcs $\gamma_{1}$ and $\gamma_{2}$ over the same parameter interval $[a,b]$ are said to be homotopic in $\Omega$ if there exists a continuous function $\gamma(t,u)$, defined on a rectangle $[a,b] \times [0,1]$, with the following properties:
1. $\gamma(t,u) \in \Omega$ for all $(t,u)$.

2. $\gamma(t,0) = \gamma_1(t)$, $\gamma(t,1) = \gamma_2(t)$ for all $t$.

3. $\gamma(a,u) = \gamma_1(a) = \gamma_2(a)$, $\gamma(b,u) = \gamma_1(b) = \gamma_2(b)$ for all $u$.
\end{originalstatement}
It is only for the sake of convenience that we have required the parametric intervals of $\gamma_{1}$ and $\gamma_{2}$ to be the same. If this is not the case, we transform the intervals into each other by a linear change of parameter, and agree to consider the original arcs as homotopic if they are homotopic in the new parametrization.

Simple formal proofs which the reader can easily supply show that the relation of homotopy, as defined above, is an equivalence relation. We can thus divide all arcs into equivalence classes, called homotopy classes; the arcs in a homotopy class have common end points and can be deformed into each other within $\Omega$. It deserves to be pointed out that different parametric representations of the same arc are always homotopic. Indeed, $\gamma_{2}(t)$ is a reparametrization of $\gamma_{1}(t)$ if and only if there exists a nondecreasing function $\tau(t)$ such that $\gamma_{2}(t) = \gamma_{1}(\tau(t))$. The function $\gamma(t,u) =$
\bookfigure{308}{1}{FIG. 8-2. Deformation.}
% Original PDF physical page 309
\sourcepage{309}
$\gamma_{1}((1 - u)t + u\tau(t))$ has all its values on the arc under consideration, and therefore in $\Omega$. For $u = 0$ and $u = 1$ we obtain respectively $\gamma(t,0) = \gamma_{1}(t)$ and $\gamma(t,1) = \gamma_{1}(\tau(t)) = \gamma_{2}(t)$ as required, and the end points are fixed.

If two arcs $\gamma_{1}$ and $\gamma_{2}$ are traced in succession, with $\gamma_{2}$ beginning at the terminal point of $\gamma_{1}$, they form a new arc which we will now denote by $\gamma_{1}\gamma_{2}$ in contrast to the notation $\gamma_{1} + \gamma_{2}$ preferred in homology theory. The parametrization of $\gamma_{1}\gamma_{2}$ is not uniquely determined, but for the determination of the homotopy\index{Homotopy} class this is of no importance. Very simple reasoning shows, moreover, that the homotopy class of $\gamma_{1}\gamma_{2}$ depends only on the homotopy classes of $\gamma_{1}$ and $\gamma_{2}$. By virtue of this fundamental fact we may consider the operation which leads to the homotopy class of $\gamma_{1}\gamma_{2}$ as a multiplication of homotopy classes. It is defined only when the initial point of $\gamma_{2}$ coincides with the terminal point of $\gamma_{1}$. If we restrict our attention to the homotopy classes of closed curves which begin and end at a fixed point $z_{0}$, the product is always defined and is represented by a curve in the same family. What is more, with this definition of product the homotopy classes of closed curves from $z_{0}$, with respect to the region $\Omega$, form a group. In order to prove this assertion we must establish:

1. The associative law: $(\gamma_{1}\gamma_{2})\gamma_{3}$ is homotopic to $\gamma_{1}(\gamma_{2}\gamma_{3})$.

2. Existence of a unit curve\index{Curve!unit}\index{Unit curve} 1 such that $\gamma 1$ and $1\gamma$ are homotopic to $\gamma$.

3. Existence of an inverse $\gamma^{-1}$ such that $\gamma\gamma^{-1}$ and $\gamma^{-1}\gamma$ are homotopic to 1.

The associative law is trivial since $(\gamma_1\gamma_2)\gamma_3$ is at most a reparametrization of $\gamma_{1}(\gamma_{2}\gamma_{3})$. For a unit curve we can choose the constant $z = z_0$; actually, the symbol 1 may represent any closed curve which can be shrunk to the point $z_0$. Finally, the inverse $\gamma^{-1}$ is the curve $\gamma$ traced in the opposite direction. If $\gamma$ is represented by $z = \gamma(t)$, $t \in [a,b]$, $\gamma^{-1}$ can be represented by $z = \gamma(2b - t)$, $t \in [b, 2b - a]$. The equation of $\gamma\gamma^{-1}$ is thus
\begin{equation*}
z = \left\{ \begin{array}{l l} \gamma (t) & \text {   for   } a \leq t \leq b \\ \gamma (2 b - t) & \text {   for   } b \leq t \leq 2 b - a. \end{array} \right.
\end{equation*}
The curve can be shrunk to a point by means of the deformation
\begin{equation*}
\gamma(t,u)=\left\{\begin{array}{@{}l@{\;}l@{}}
\gamma(t)&\text{for }a\leq t\leq ua+(1-u)b,\\
\gamma(ua+(1-u)b)&\text{for }ua+(1-u)b\leq t\leq u(b-a)+b,\\
\gamma(2b-t)&\text{for }u(b-a)+b\leq t\leq2b-a.
\end{array}\right.
\end{equation*}
The interpretation is clear: we are letting the turning point recede from $\gamma(b)$ to $\gamma(a)$. Since $\gamma(t,1) = \gamma(a) = z_0$ for all $t \in [a, 2b - a]$ we have proved that $\gamma\gamma^{-1}$ is homotopic to 1. The proof is independent of the hypothesis that $\gamma$ be a closed curve; thus $\gamma\gamma^{-1}$ is homotopic to 1 for any arc $\gamma$ from $z_0$.


% Original PDF physical page 310
\sourcepage{310}
The group which we have constructed is called the homotopy\index{Homotopy} group, or the fundamental group\index{Fundamental group}, of the region $\Omega$ with respect to the point $z_0$. As an abstract group it does not depend on the point $z_0$. If $z_0'$ is another point in $\Omega$, we join $z_0$ to $z_0'$ by an arc $c$ in $\Omega$. To every closed curve $\gamma'$ from $z_0'$ corresponds a closed curve $\gamma = c\gamma'c^{-1}$ from $z_0$. This correspondence is homotopy preserving and may thus be regarded as a correspondence between homotopy classes. As such it is product preserving, for $(c\gamma_1'c^{-1})(c\gamma_2'c^{-1})$ is homotopic to $c(\gamma_1'\gamma_2')c^{-1}$, by cancellation of $c^{-1}c$. Finally, the correspondence is one to one, for if $\gamma$ is given we can choose $\gamma' = c^{-1}\gamma c$ and find that the corresponding curve $c\gamma'c^{-1} = (cc^{-1})\gamma(cc^{-1})$ is homotopic to $\gamma$. It is thus proved that the homotopy groups with respect to $z_0$ and $z_0'$ are isomorphic.

If $\gamma_{1}, \gamma_{2}$ are any two arcs with the initial point $z_{0}$ and a common terminal point, then $\gamma_{1}$ is homotopic to $\gamma_{2}$ if and only if $\gamma_{1}\gamma_{2}^{-1}$ is homotopic to 1. For if $\gamma_{1}$ is homotopic to $\gamma_{2}$, then $\gamma_{1}\gamma_{2}^{-1}$ is homotopic to $\gamma_{2}\gamma_{2}^{-1}$, and hence to 1. Conversely, if $\gamma_{1}\gamma_{2}^{-1}$ is homotopic to 1, then
\begin{equation*}
(\gamma_ {1} \gamma_ {2} ^ {- 1}) \gamma_ {2} = \gamma_ {1} (\gamma_ {2} ^ {- 1} \gamma_ {2})
\end{equation*}
is simultaneously homotopic to $\gamma_{1}$ and $\gamma_{2}$, proving that $\gamma_{1}$ is homotopic to $\gamma_{2}$. For this reason it is sufficient to study the homotopy of closed curves.

The explicit determination of homotopy groups is simplified by the fact that the homotopy group is obviously a topological invariant. Indeed, by a topological mapping of $\Omega$ onto $\Omega'$ any deformation in $\Omega$ can be carried over to $\Omega'$ and is seen to determine a product preserving one-to-one correspondence between the homotopy classes. Topologically equivalent regions have therefore isomorphic homotopy groups.

The homotopy group of a disk reduces to the unit element; this means that any two arcs with common end points are homotopic. The proof makes use of the convexity of the disk: the arc $z = \gamma_{1}(t)$ can be deformed into $z = \gamma_{2}(t)$ by means of the deformation
\begin{equation*}
\gamma (t, u) = (1 - u) \gamma_ {1} (t) + u \gamma_ {2} (t)
\end{equation*}
whose deformation paths are line segments. The same proof would be valid for any convex region. In particular, the whole plane has likewise a homotopy group which reduces to the unit element.

We proved in Chap.~\hyperref[chap:6]{6}, Sec.~\hyperref[sec:6.1]{1}, that any simply connected region which is not the whole plane can be mapped conformally onto a disk. In this connection the conformality is not important, but the fact that the mapping is topological permits us to conclude that any simply connected region has a fundamental group which reduces to its unit element. We shall find that the converse is also true.


% Original PDF physical page 311
\sourcepage{311}
\setcounter{section}{1}\setcounter{subsection}{5}\subsection{The Monodromy Theorem}\label{sec:8.1.6}
Let $\Omega$ be a fixed region in the complex plane. We consider the case of a global analytic\index{Function!analytic!global} function\index{Global analytic function} $\mathbf{f}$ which can be continued along all arcs $\gamma$ contained in $\Omega$, starting from any of its germs defined at the initial point $\zeta_0$ of $\gamma$. More precisely, for any function element $(f_0,\Omega_0)$ of $\mathbf{f}$ with $\zeta_0\in \Omega_0$, there shall exist a continuation $\overline{\gamma}$ over $\gamma$ beginning with the germ defined by $(f_0,\zeta_0)$.

When two arcs $\gamma_{1}$, $\gamma_{2}$ with common end points are given, we are interested to know whether a common initial germ, continued along $\gamma_{1}$ and $\gamma_{2}$, will lead to the same germ over the terminal point. The basic theorem, known as the monodromy theorem\index{Monodromy theorem|(}\index{Monodromy theorem}, is the following:

\begin{originalstatement}{Theorem}{2}\label{thm:8.2}
If the arcs $\gamma_{1}$ and $\gamma_{2}$ are homotopic in $\Omega$, and if a given germ of $\mathbf{f}$ at the initial point can be continued along all arcs contained in $\Omega$, then the continuations of this germ along $\gamma_{1}$ and $\gamma_{2}$ lead to the same germ at the terminal point.
\end{originalstatement}
To begin with we note that continuation along an arc of the form $\gamma\gamma^{-1}$ will always lead back to the initial germ. Similarly, continuation along an arc of the form $\sigma_{1}(\gamma\gamma^{-1})\sigma_{2}$ will have the same effect as continuation along $\sigma_{1}\sigma_{2}$. For this reason, to say that the continuations along $\gamma_{1}$ and $\gamma_{2}$ lead to the same end result is equivalent to saying that continuation along $\gamma_{1}\gamma_{2}^{-1}$ leads back to the initial germ.

According to the assumption there exists a deformation $\gamma(t,u)$ of $\gamma_{1}$ into $\gamma_{2}$. Every arc $\sigma$ in the deformation rectangle $R = [a,b] \times [0,1]$ is carried by $\gamma(t,u)$ into an arc $\sigma' \in \Omega$, and if $\sigma'$ begins at the initial point of $\gamma_{1}$ and $\gamma_{2}$, there exists a unique continuation along $\sigma'$ from the initial germ; for simplicity we shall call it a continuation along $\sigma$. The theorem asserts that the continuation along the perimeter $\Gamma$ of $R$ leads back to the initial germ. The sense in which $\Gamma$ is described is immaterial, but should be fixed once and for all.

A simple proof can be based on the method of bisection. We begin by bisecting R horizontally, and denote by $\pi_{1}$ the perimeter of the lower half $R_{1}$, described from the lower left-hand corner 0 and in the direction which coincides with the direction of $\Gamma$ along the common side. With the upper half $R_{2}$ we associate a curve $\pi_{2}$ which begins at 0, leads vertically to the lower left-hand corner of $R_{2}$, describes the perimeter of $R_{2}$ in the sense which coincides with that of $\Gamma$ along the common side, and returns vertically to 0 (Fig.~\hyperref[fig:8-3]{8-3}). We recognize that the curve $\pi_{1}\pi_{2}$ differs from $\Gamma$ only by an intermediate arc of the form $\sigma\sigma^{-1}$. For this reason the effect of continuing along $\pi_{1}\pi_{2}$ is the same as if we continue along $\Gamma$. Consequently, if $\pi_{1}$ and $\pi_{2}$ both lead back to the initial germ, so does $\Gamma$.


% Original PDF physical page 312
\sourcepage{312}
\bookfigure{312}{1}{FIG. 8-3. The monodromy theorem.}
We make now the opposite assumption that $\Gamma$ does not lead back to the initial germ. Then either $\pi_1$ or $\pi_2$ has the same property. The corresponding rectangle is bisected vertically, and the same reasoning is applied. When the process is repeated, we obtain a sequence of rectangles $R\supset R^{(1)}\supset R^{(2)}\supset \dots \supset R^{(n)}\supset \dots$ and corresponding closed curves $\pi^{(n)}$ such that the continuation of the initial germ along $\pi^{(n)}$ does not lead back to the same germ. Each $\pi^{(n)}$ is of the form $\sigma_n\Gamma_n\sigma_n^{-1}$ where $\sigma_n$ is a well-determined polygon leading from 0 to the lower left-hand corner of $R^{(n)}$ and $\Gamma_n$ denotes the perimeter of $R^{(n)}$; moreover, $\sigma_n$ is a subarc of $\sigma_{n+1}$.

As $n \to \infty$ the rectangles $R^{(n)}$ converge to a point $P_{\infty}$, and the polygons $\sigma_{n}$ form, in the limit, an arc $\sigma_{\infty}$ ending at $P_{\infty}$. There exists a continuation of the initial germ along $\sigma_{\infty}$; it terminates with a germ determined by a function element $(f_{\infty}, \Omega_{\infty})$ over the image $\zeta_{\infty}$ of $P_{\infty}$ under the mapping $\gamma(t, u)$. For sufficiently large n the image of $\Gamma_{n}$ will be contained in $\Omega_{\infty}$, and the germ obtained at the terminal point of $\sigma_{n}$ will belong to the function element $(f_{\infty}, \Omega_{\infty})$. When this is the case, the element $(f_{\infty}, \Omega_{\infty})$ can be used to construct a continuation along $\pi^{(n)}$ which leads back to the initial germ. This contradicts the property by which $\pi^{(n)}$ was chosen, and we have proved that the continuation along $\Gamma$ ends with the initial germ.

The monodromy theorem\index{Monodromy theorem} implies, above all, that any global analytic\index{Function!analytic!global} function\index{Global analytic function} which can be continued along all arcs in a simply connected region determines one single-valued analytic function for each choice of the initial branch. This fact can also be expressed by saying that a Riemann surface (without branch points) over a simply connected region must consist of a single sheet.

We can further draw the consequence, already announced, that a region whose homotopy\index{Homotopy} group reduces to the unit element must necessarily be simply connected. For suppose that $\Omega$ is multiply connected. Then there exists a bounded component $E_{0}$ of the complement of $\Omega$, and
% Original PDF physical page 313
\sourcepage{313}
if $z_{0} \in E_{0}$ we know that $\log (z - z_{0})$ is not single-valued in $\Omega$. By the monodromy theorem\index{Monodromy theorem} it follows that the homotopy\index{Homotopy} group of $\Omega$ cannot reduce to the unit element.

This is the last step toward proving the equivalence of the following three characterizations of simply connected regions: (1) $\Omega$ is simply connected if its complement is connected; (2) $\Omega$ is simply connected if it is homeomorphic with a disk; (3) $\Omega$ is simply connected if its fundamental group reduces to the unit element.

\setcounter{section}{1}\setcounter{subsection}{6}\subsection{Branch Points}\label{sec:8.1.7}
For a closer study of the singularities of multiple-valued functions it is necessary to determine, explicitly, the fundamental group of a punctured disk. Let the punctured disk be represented by $0 < |z| < \rho$, and consider a fixed point, for instance the point $z_{0} = r < \rho$ on the positive radius. By means of a central projection, given by
\begin{equation*}
\gamma (t, u) = (1 - u) \gamma (t) + u r \frac {\gamma (t)}{| \gamma (t) |},
\end{equation*}
any closed curve $\gamma$ from $z_{0}$ can be deformed into a curve which lies on the circle $|z| = r$. It is thus sufficient to consider curves on that circle. We continue to use the notation $\gamma(t)$.

By continuity every $t_0$ has a neighborhood in which $|\gamma(t) - \gamma(t_0)| < r$; in such a neighborhood $\gamma(t)$ cannot take both the values $r$ and $-r$. It follows easily, by use of Heine-Borel's lemma or the method of bisection, that it is possible to write $\gamma = \gamma_1\gamma_2 \cdots \gamma_n$ where each $\gamma_k$ either does not pass through $r$ or does not pass through $-r$. For simplicity, let us refer to the points $r$ and $-r$ by letters $P_0$ and $P_0'$ (Fig.~\hyperref[fig:8-4]{8-4}), and let the end points of $\gamma_k$ be denoted by $P_k$ and $P_{k+1}$. Since $\gamma_k$ is contained in the simply connected region obtained by deleting either the positive or negative radius, it can be deformed into one of the two arcs $P_kP_{k+1}$. As a result $\gamma$ can be deformed into a product of simple arcs with the successive end points $P_0P_1P_2 \cdots P_nP_0$. This path may in turn be replaced by $P_0P_1P_2P_0P_2P_3P_0 \cdots P_0P_{n-1}P_nP_0$ where each arc $P_kP_0$ and $P_0P_k$ is, for definiteness, the one which does not contain $P_0'$. In fact, the new path is obtained by inserting the doubly traced arcs $P_kP_0P_k$ which we know to be homotopic to 1.

We have shown that each $\gamma$ is homotopic to a product of closed curves of the form $P_{0}P_{k}P_{k+1}P_{0}$. If $P_{k}P_{k+1}$ does not contain $P_{0}^{\prime}$, this curve is homotopic to 1. If, on the other hand, $P_{k}P_{k+1}$ contains $P_{0}^{\prime}$ it is seen by enumeration of the possible cases that the curve is homotopic to C or $C^{-1}$, where C is the full circle. Consequently, every closed curve is homotopic to a power of C.


% Original PDF physical page 314
\index{Monodromy theorem|)}\sourcepage{314}
\bookfigure{314}{1}{FIG. 8-4}
Finally, we observe that $C^m$ is homotopic to 1 only if $m = 0$. This is seen by the fact that
\begin{equation*}
\int_ {C ^ {m}} \frac {d z}{z} = m \cdot 2 \pi i,
\end{equation*}
while if the curve were homotopic to 1 the integral would have to vanish. From our results we conclude that the fundamental group of the punctured disk is isomorphic to the additive group\index{Additive group} of integers. Evidently, an arbitrary annulus has the same fundamental group.

We consider now a global analytic\index{Function!analytic!global} function\index{Global analytic function} $\mathbf{f}$ which can be continued along all arcs in the punctured disk $0 < |z| < \rho$. We choose an initial germ at $z_{0} = r$ and continue it along all curves $C^{m}$. Either the continuation never returns to the initial germ, or there exists a smallest positive integer h such that $C^{h}$ leads back to where we started. In the latter case, set $m = nh + q$ with n an integer and $0 \leq q < h$. If $C^{m}$ leads back to the initial germ, so does $C^{q}$. Because of the choice of m this is possible only if q = 0. Thus $C^{m}$ leads to the initial germ only if m is a multiple of h.

Consider the mapping $z = \zeta^h$ of $0 < |\zeta| < \rho^{1/h}$ on $0 < |z| < \rho$. We claim that $\mathbf{f}$ can be expressed as a single-valued analytic function $F(\zeta)$ in the following sense: For every $\zeta_1, 0 < |\zeta_1| < \rho^{1/h}$, there exists a function element $(f,\Omega) \in \mathbf{f}$ with $\zeta_1^h \in \Omega$, such that $F(\zeta) = f(\zeta^h)$ in a neighborhood of $\zeta_1$; in particular, it is required that $\zeta_0 = r^{1/h}$ corresponds in this way to the initial germ of $\mathbf{f}$ at $z_0$.

In order to construct $F(\zeta)$ we join $\zeta_0$ to $\zeta$ by an arc $\gamma'$ and continue the initial germ of $\mathbf{f}$ along the image of $\gamma'$ under the mapping $z = \zeta^h$; we define $F(\zeta)$ to be the value of the terminal germ under this continuation. It must be proved that $F(\zeta)$ is uniquely determined. If $\zeta_1'$ and $\zeta_2'$ are two paths
% Original PDF physical page 315
\sourcepage{315}
from $\zeta_0$ to $\zeta$, then $\zeta_1'\zeta_2'^{-1}$ can be deformed into a power $C''^n$ of the circle through $\zeta_0$. Consequently, the image curve $\zeta_1\zeta_2^{-1}$ can be deformed into the image of $C''^n$, which is $C^{nh}$. But $C^{nh}$ leads back to the initial germ, and hence $\zeta_1$ and $\zeta_2$ determine the same value $F(\zeta)$. Finally, if $\zeta$ is in a neighborhood of $\zeta_1$, we can first follow an arc $\zeta_1'$ from $\zeta_0$ to $\zeta_1$ and then a variable arc $\gamma'$ from $\zeta_1$ to $\zeta$ which stays within the neighborhood. If the neighborhood is sufficiently restricted, the continuation along the image of $\gamma'$ is determined by a single function element $(f,\Omega)$, and $F(\zeta) = f(\zeta^h)$ in that neighborhood.

Since $F(\zeta)$ is single-valued and analytic in a punctured neighborhood of the origin, it has a convergent Laurent development of the form
\begin{equation*}
F (\zeta) = \sum_ {- \infty} ^ {\infty} A _ {n} \zeta^ {n}.\tag{1}\label{eq:8.1}
\end{equation*}
It must be observed that this development depends on the choice of the initial germ; different choices may yield entirely different developments and, in particular, different values of h. Actually, even the series \hyperref[eq:8.1]{(1)} yields h different developments, corresponding to the h initial values of $z^{1/h}$. If we write $\omega = e^{2\pi i/h}$, these developments are represented by
\begin{equation*}
f _ {\nu} (z) = \sum_ {- \infty} ^ {\infty} A _ {n} \omega^ {\nu n} z ^ {n / h} \quad (\nu = 0, 1, \dots , h - 1).\tag{2}\label{eq:8.2}
\end{equation*}
When the germ $(f_{\nu},z_{0})$ is continued along C it leads to $(f_{\nu+1},z_{0})$, with the understanding that the subscript h is identified with 0.

In special cases the Laurent development may contain only a finite number of negative powers. Then $F(\zeta)$ has either a removable singularity or a pole, and the multiple-valued function $f(z)$ (or, more correctly, the global analytic\index{Function!analytic!global} function\index{Global analytic function} obtained by continuing the given initial branch within a punctured disk) is said to have an algebraic singularity\index{Algebraic singularity} or branch point\index{Branch point}\index{Point!branch} at z = 0, provided of course that h > 1. If $F(\zeta)$ has a removable singularity, the branch point is an ordinary algebraic singularity, in the opposite case it is an algebraic pole\index{Pole!algebraic}. In either case $f(z)$ tends to a definite limit $A_{0}$ or $\infty$ as z tends to 0 along an arbitrary arc.

Clearly, we could just as well have studied an isolated singularity at an arbitrary point a or $\infty$, and the radius of the punctured disk can be as small as we wish. In the case of a finite h the correspondence between $w = f(z)$ and the independent variable z can be expressed through equations of the form
\begin{equation*}
\begin{array}{l} w = \sum_ {- \infty} ^ {\infty} A _ {n} \zeta^ {n} \\ z = a + \zeta^ {h} \quad \text { or } \quad z = \zeta^ {- h}. \end{array}
\end{equation*}
% Original PDF physical page 316
\sourcepage{316}
The\index{Irreducible polynomial|(} variable $\zeta$ takes the name of \emph{local uniformizing variable}\index{Uniformizing variable}.

In the case of an algebraic singularity it is desirable to complete the Riemann surface of f so as to include a branch point with the projection a. The branch point itself is not a germ of f, but it is fully determined by a set of fractional power series developments
\begin{equation*}
f _ {\nu} (z) = \sum_ {\nu = \nu_ {0}} ^ {\infty} A _ {n} \omega^ {\nu n} (z - a) ^ {n / h}\tag{3}\label{eq:8.3}
\end{equation*}
analogous to \hyperref[eq:8.2]{(2)}; for a singularity at infinity z - a has to be replaced by 1/z. The neighborhoods of the branch point shall include the branch point itself as well as, for some $\delta > 0$, all germs $(f_{\nu}, \zeta)$ with $|\zeta - a| < \delta$ obtained by substituting in \hyperref[eq:8.3]{(3)} a single-valued branch of $(z - a)^{1/h}$, defined in a neighborhood of $\zeta$. The resulting topological space will be a \emph{surface} in the sense that every point, including the branch points, has a neighborhood which is homeomorphic to a disk.

In the Weierstrass theory it is customary to consider the totality of all power series developments, including the fractional ones, that are obtainable by analytic continuation from a single one, and to call it an \emph{analytic configuration} (analytisches Gebilde).

\setcounter{section}{1}\section{Algebraic Functions}\label{sec:8.2}
An equation of the form $P(w,z)=0$, where P is a polynomial in two variables, has for each z a finite number of solutions $w_{1}(z),\ldots,w_{m}(z)$. We wish to show that these roots can be interpreted as values of a global analytic\index{Function!analytic!global} function\index{Global analytic function} $\mathbf{f}(z)$ which is then called an \emph{algebraic function}. Conversely, if a global analytic function is given, we want to be able to tell whether it does or does not satisfy a polynomial equation.

\setcounter{section}{2}\setcounter{subsection}{0}\subsection{The Resultant of Two Polynomials}\label{sec:8.2.1}
A polynomial $P(w,z)$ in two variables is \emph{irreducible} if it cannot be expressed as the product of two polynomials none of which is constant. Two polynomials P and Q are \emph{relatively prime} if they have no common factor except for constants.

The following theorem is algebraic in character. Because of its fundamental importance for the theory of algebraic functions\index{Algebraic function|(}\index{Function!algebraic|(}\index{Algebraic function}\index{Function!algebraic} we will nevertheless reproduce its proof.

\begin{originalstatement}{Theorem}{3}\label{thm:8.3}
If $P(w,z)$ and $Q(w,z)$ are relatively prime\index{Relatively prime} polynomials, there are only a finite number of values $z_0$ for which the equations $P(w,z_0) = 0$ and $Q(w,z_0) = 0$ have a common root.
\end{originalstatement}

% Original PDF physical page 317
\index{Homotopy|)}\index{Irreducible polynomial|)}\sourcepage{317}
We suppose that P and Q are ordered according to decreasing powers of w and set $Q(w,z) = b_{0}(z)w^{m} + \cdots + b_{m}(z)$ where $b_{0}(z)$ is not identically zero. If P is divided by Q, the division algorithm yields a quotient and remainder which are polynomials in w and rational functions in z. We set up a Euclidean algorithm of the form
\begin{equation*}
\begin{array}{l} c _ {0} P = q _ {0} Q + R _ {1} \\ c _ {1} Q = q _ {1} R _ {1} + R _ {2} \\ c _ {2} R _ {1} = q _ {2} R _ {2} + R _ {3} \\ \dots \dots \\ c _ {n - 1} R _ {n - 2} = q _ {n - 1} R _ {n - 1} + R _ {n} \end{array}\tag{4}\label{eq:8.4}
\end{equation*}
where the $q_{k}$ and $R_{k}$ are polynomials in w and z while the $c_{k}$ are polynomials in z used to clear the fractions. The degrees in w of the $R_{k}$ are decreasing, and $R_{n}$ is a polynomial in z alone. If $R_{n}(z)$ were identically zero, the unique factorization theorem implies, by the last relation in \hyperref[eq:8.4]{(4)}, that $R_{n-2}$ would be divisible by any irreducible factor of $R_{n-1}$ which is of positive degree in w. The same reasoning shows, step by step, that all the $R_{k}$ as well as Q and P would be divisible by the same factor. This is contrary to the assumption, for $R_{n-1}$ is of positive degree in w and must therefore have an irreducible factor which contains w.

Suppose now that $P(w_0, z_0) = 0$ and $Q(w_0, z_0) = 0$. Substituting these values in \hyperref[eq:8.4]{(4)} we obtain $R_1(w_0, z_0) = 0, \ldots, R_{n-1}(w_0, z_0) = 0$ and finally $R_n(z_0) = 0$. But since $R_n$ is not identically zero, there are only a finite number of $z_0$ which satisfy this condition, and the theorem follows.

The polynomial $R_{n}(z)$ is called the \emph{resultant} of P and Q. More precisely, if we wish the resultant\index{Resultant} to be uniquely determined, we should require that the exponents $c_{k}$ in \hyperref[eq:8.3]{(3)} are of the lowest degree possible. We are not so interested in the resultant as in the statement of Theorem~\hyperref[thm:8.3]{3}. The theorem will be applied to an irreducible polynomial $P(w,z)$ and its partial derivative $P_{w}(w,z)$ with respect to w. These polynomials are relatively prime as soon as P has positive degree in w, and the resultant of P and $P_{w}$ is called the \emph{discriminant} of P. The zeros of the discriminant\index{Discriminant} are the values $z_{0}$ for which the equation $P(w,z_{0}) = 0$ has multiple roots.

We note, finally, that the resultant $R(z)$ of any two relatively prime polynomials P and Q can be written in the form $R = pP + qQ$ where p and q are polynomials. This follows immediately from \hyperref[eq:8.4]{(4)}.

\setcounter{section}{2}\setcounter{subsection}{1}\subsection{Definition and Properties of Algebraic Functions}\label{sec:8.2.2}
We begin by formulating a precise definition:

\begin{originalstatement}{Definition}{3}\label{def:8.3}
A global analytic\index{Function!analytic!global} function\index{Global analytic function} $\mathbf{f}$ is called an algebraic function\index{Algebraic function}\index{Function!algebraic} if all its function elements $(f,\Omega)$ satisfy a relation $P(f(z),z) = 0$ in $\Omega$,
% Original PDF physical page 318
\sourcepage{318}
where $P(w,z)$ is a polynomial which does not vanish identically.

\end{originalstatement}
Because of the permanence of functional relations it is sufficient to assume that one function element satisfies the equation $P(f(z),z) = 0$. The others will then automatically satisfy the same relation. Moreover, it may be assumed that $P(w,z)$ is an irreducible polynomial. Suppose indeed that $P(w,z)$ has the factorization $P = P_1P_2\ldots P_n$ in irreducible factors. For any fixed point $z \in \Omega$ one of the equations $P_k(f(z),z) = 0$ must hold. If we consider a sequence of different points $z_n \in \Omega$ which tend to a limit in $\Omega$, then one of the relations $P_k(f(z_n),z_n)=0$ must hold infinitely often. It follows that this particular relation $P_k(f(z),z) = 0$ is satisfied identically in $\Omega$ and, consequently, by all the function elements of $f$. We are thus free to replace $P$ by $P_k$.

It is also easy to see that the irreducible polynomial P determined by an algebraic function\index{Algebraic function}\index{Function!algebraic} is unique up to a constant factor. If Q is an essentially different irreducible polynomial, we can determine the resultant $R(z) = pP + qQ$. If $P(f(z), z) = 0$ and $Q(f(z), z) = 0$ for all $z \in \Omega$ we would obtain $R(z) = 0$ in $\Omega$, contrary to the fact that $R(z)$ is not identically zero. We note that P cannot reduce to a polynomial of z alone. If it contains only w, it must be of the form w - a, and the function f reduces to the constant a.

We prove next that there exists an algebraic function corresponding to any irreducible polynomial $P(w,z)$ of positive degree in w. Suppose that
\begin{equation*}
P (w, z) = a _ {0} (z) w ^ {n} + a _ {1} (z) w ^ {n - 1} + \dots + a _ {n} (z).
\end{equation*}
If $z_{0}$ is neither a zero of the polynomial $a_{0}(z)$ nor a zero of the discriminant of P, the equation $P(w,z_{0}) = 0$ has exactly n distinct roots $w_{1}, w_{2}, \ldots, w_{n}$. Under this condition the following is true:

\begin{originalstatement}{Lemma}{1}\label{lem:8.1}
There exists an open disk $\Delta$, containing $z_0$, and $n$ function elements $(f_1, \Delta)$, $(f_2, \Delta)$, $\ldots$, $(f_n, \Delta)$ with these properties:
(a) $P(f_{i}(z),z) = 0$ in $\Delta$;

(b) $f_{i}(z_{0}) = w_{i}$;

(c) if $P(w,z) = 0$, $z \in \Delta$, then $w = f_i(z)$ for some $i$.

\end{originalstatement}
The polynomial $P(w, z_0)$ has simple zeros at $w = w_i$. We determine $\varepsilon > 0$ so that the disks $|w - w_i| \leq \varepsilon$ do not overlap and denote the circles $|w - w_i| = \varepsilon$ by $C_i$. Then $P(w, z_0) \neq 0$ on $C_i$, and by the argument principle
\begin{equation*}
\frac {1}{2 \pi i} \int_ {C _ {i}} \frac {P _ {w} (w , z _ {0})}{P (w , z _ {0})} d w = 1.
\end{equation*}
% Original PDF physical page 319
\sourcepage{319}
If $z_{0}$ is replaced by z, the integrals become well-defined continuous functions of z in a neighborhood of $z_{0}$. Since they can only take integer values, there exists a neighborhood $\Delta$ such that
\begin{equation*}
\frac {1}{2 \pi i} \int_ {C _ {i}} \frac {P _ {w} (w , z)}{P (w , z)} d w = 1
\end{equation*}
for all $z \in \Delta$. This means that the equation $P(w, z) = 0$ has exactly one root in the disk $|w - w_i| < \varepsilon$; we denote this root by $f_i(z)$. By the residue calculus its value is given by
\begin{equation*}
f _ {i} (z) = \frac {1}{2 \pi i} \int_ {C _ {i}} w \frac {P _ {w} (w , z)}{P (w , z)} d w.
\end{equation*}
This representation shows that $f_{i}(z)$ is analytic. Moreover, $f_{i}(z_{0}) = w_{i}$, and (c) follows from the fact that we have exhibited n roots of the equation $P(w,z) = 0$, and it can have no more.

The lemma implies at once that there exists an algebraic function\index{Algebraic function}\index{Function!algebraic} f corresponding to the polynomial P; in fact, we can choose f to be the global analytic\index{Function!analytic!global} function\index{Global analytic function} determined by the element $(f_{1},\Delta)$ for any $z_{0}$ which does not coincide with one of the finitely many excluded points. We will show, moreover, that all such function elements belong to the same global analytic function; this will also prove that the function f that corresponds to P is unique. Let $(f,\Omega)$ be one of these function elements. There must exist a $z_{0}\in\Omega$ which is not one of the excluded points; we determine a corresponding $\Delta$. Since $P(f(z),z)=0$ for $z\in\Omega$ it follows by (c) that $f(z)$ equals some $f_{i}(z)$ at each point of $\Delta\cap\Omega$. But then $f(z)$ equals the same $f_{i}(z)$ infinitely many times in any neighborhood of $z_{0}$, and hence $(f,\Omega)$ belongs to the same global analytic function as $(f_{i},\Delta)$.

Let the excluded points be denoted by $c_{1}, c_{2}, \ldots, c_{m}$. We wish to show that a function element $(f, \Omega)$ which satisfies $P(f(z), z) = 0$ can be continued along any arc which does not pass through a point $c_{k}$. If this were not so, there would exist an arc $\gamma[a, b]$ such that a given initial germ can be continued along all subarcs $\gamma[a, \tau]$ with $\tau < b$, but not along the whole arc. We set $z_{0} = \gamma(b)$, determine $\Delta$ according to Lemma~\hyperref[lem:8.1]{1}, and choose $\tau$ so that $\gamma(t) \in \Delta$ for all $t \in [\tau, b]$. The same reasoning as above shows that the germ $\bar{\gamma}(\tau)$ obtained by continuation along $\gamma[a, \tau]$ must be determined by one of the function elements $(f_{i}, \Delta)$. But then it can be continued all the way to b, and we have reached a contradiction.

It has not yet been proved that all elements $(f_{i},\Omega)$ belong to the same global analytic function. For this part of the proof it is necessary to study the behavior at the critical points $c_{k}$ in greater detail.


% Original PDF physical page 320
\sourcepage{320}
\setcounter{section}{2}\setcounter{subsection}{2}\subsection{Behavior at the Critical Points}\label{sec:8.2.3}
The points $c_{k}$ which so far have been excluded from our considerations were the zeros of the first coefficient $a_{0}(z)$ of P, and the zeros of the discriminant. Let $\delta$ be chosen so that the disk $|z - c_{k}| < \delta$ contains no other critical points\index{Critical point|(}\index{Critical point}. We fix a point $z_{0} \neq c_{k}$ in this disk and select one of the germs $(f_{i}, z_{0})$. This germ can be continued along all arcs in the punctured disk. Moreover, if continued along the circle C of center $c_{k}$ through $z_{0}$, it leads to a germ $(f_{j}, z_{0})$. Since there are only a finite number of choices, there must exist a smallest positive $h \leq n$ with the property that continuation along $C^{h}$ leads back to the initial germ $(f_{i}, z_{0})$. By the result of Sec.~\hyperref[sec:8.1.6]{1.6} we can write
\begin{equation*}
f _ {i} (z) = \sum_ {\nu = - \infty} ^ {\infty} A _ {\nu} (z - c _ {k}) ^ {\nu / h}.\tag{5}\label{eq:8.5}
\end{equation*}
Suppose first that $c_{k}$ is not a zero of $a_{0}(z)$. Then $f_{i}(z)$ remains bounded as z tends to $c_{k}$. Indeed, as soon as $f_{i}(z) \neq 0$ the equation $P(f_{i}(z), z) = 0$ can be written in the form
\begin{equation*}
a _ {0} (z) + a _ {1} (z) f _ {i} (z) ^ {- 1} + \dots + a _ {n} (z) f _ {i} (z) ^ {- n} = 0.\tag{6}\label{eq:8.6}
\end{equation*}
If $f_{i}(z)$ were unbounded, there would exist points $z_{n} \to c_{k}$ with $f_{i}(z_{n}) \to \infty$. Substitution in \hyperref[eq:8.6]{(6)} would yield $a_0(z_n) \to 0$, contrary to the assumption $a_0(c_k) \neq 0$. It follows that the development \hyperref[eq:8.5]{(5)} contains only positive powers, and $f_{i}$ has at most an ordinary algebraic singularity at $c_k$.

We consider now the case where $a_{0}(c_{k}) = 0$. If the multiplicity of the zero is denoted by m, we know that $\lim_{z \to c_{k}} a_{0}(z)(z - c_{k})^{-m} \neq 0$. From \hyperref[eq:8.6]{(6)} we obtain
\begin{equation*}
\begin{aligned}
&a_0(z)(z-c_k)^{-m}+a_1(z)(z-c_k)^{-m}f_i(z)^{-1}+\dots\\
&\qquad+a_n(z)(z-c_k)^{-m}f_i(z)^{-n}=0.
\end{aligned}
\end{equation*}
If the expression $f_{i}(z)(z - c_{k})^{m}$ were unbounded, we would again be led to a contradiction. As in Sec.~\hyperref[sec:8.1.7]{1.7} we write
\begin{equation*}
F (\zeta) = \sum_ {- \infty} ^ {\infty} A _ {\nu} \zeta^ {\nu}
\end{equation*}
and find that $F(\zeta)\zeta^{mh}$ is bounded. Consequently $F(\zeta)$ has a pole of at most order mh, and $f_{i}$ has at most an algebraic pole at $c_{k}$ or, in special cases, an ordinary algebraic singularity.

Finally, the behavior at $z = \infty$ needs also to be discussed. It is clear that we have a development of the form
\begin{equation*}
f _ {i} (z) = \sum_ {- \infty} ^ {\infty} A _ {\nu} z ^ {\nu / h},
\end{equation*}
% Original PDF physical page 321
\sourcepage{321}
valid in a neighborhood of $\infty$. Suppose that the polynomial $a_{i}(z)$ is of degree $r_{i}$ (the coefficients which vanish identically will be left out of consideration). Choose an integer m such that
\begin{equation*}
m > \frac {1}{k} (r _ {k} - r _ {0})\tag{7}\label{eq:8.7}
\end{equation*}
for $k = 1, \ldots, n$. We contend that $f_{i}(z)z^{-m}$ must be bounded as $z \to \infty$. If this were not so we would have $f_{i}(z)^{-1}z^{m} \to 0$ for a sequence tending to $\infty$. This would imply $f_{i}(z)^{-k}z^{mk} \to 0$ and, by \hyperref[eq:8.7]{(7)}, $f_{i}(z)^{-k}z^{r_{k}-r_{0}} \to 0$ for $k \geq 1$. If \hyperref[eq:8.6]{(6)} is multiplied by $z^{-r_{0}}$ it follows that all terms except the first tend to zero. This is a contradiction, and we may conclude that $f_{i}(z)$ has at most an algebraic pole at infinity.

To sum up, we have proved that an algebraic function\index{Algebraic function}\index{Function!algebraic} has at most algebraic singularities in the extended plane. We will now prove a converse of this statement. In order to obtain a converse it is essential to add an assumption which implies that there are only a finite number of branches at a given point.

Let $\mathbf{f}$ be a global analytic\index{Function!analytic!global} function\index{Global analytic function}. For each $c$ we assume the existence of a punctured disk $\Delta$, centered at $c$, such that all germs of $\mathbf{f}$ which are defined at a point $z_0 \in \Delta$ can be continued along all arcs in $\Delta$ and show algebraic character at $c$. The assumption shall be satisfied also for $c = \infty$, in which case $\Delta$ is the exterior of a circle. Moreover, for one $\Delta$ it must be assumed that the number of different germs at $z_0$ is finite.

Since the extended plane can be covered by a finite number of disks $\Delta$, the center included, it follows that only a finite number of points c can be effective singularities; we denote these points by $c_{k}$. It is easy to prove that the number of germs at any point $z \neq c_{k}$ is constant. For every such point has a neighborhood in which all germs of f are single-valued and can be continued throughout the neighborhood. It follows that the set of points z with exactly n germs is open (n can be finite or infinite). Since the extended plane minus the points $c_{k}$ is connected, only one of these sets is nonempty. Hence n is constant, by assumption it cannot be infinite, and it cannot be zero since in that case f would be an empty collection of function elements.

The branches at any point $z \neq c_{k}$ may now be denoted as $f_{1}(z), \ldots, f_{n}(z)$, except that the ordering remains indeterminate. We form now the elementary symmetric functions of the $f_{i}(z)$, that is to say the coefficients of the polynomial
\begin{equation*}
(w - f _ {1} (z)) (w - f _ {2} (z)) \cdot \cdot \cdot (w - f _ {n} (z)).
\end{equation*}
These coefficients are well-defined functions of z, and obviously analytic except for possible isolated singularities at the points $c_{k}$. As z approaches
% Original PDF physical page 322
\sourcepage{322}
$c_{k}$\index{Linear differential equation|(} we know that each $f_{i}(z)$ may grow toward infinity at most like a negative power of $|z - c_{k}|$. The same is hence true of the elementary symmetric functions. We conclude that the isolated singularities, including the one at infinity, are at most poles, and consequently the elementary symmetric functions are rational functions of z. If their common denominator is denoted by $a_{0}(z)$, we find that all branches $f_{i}(z)$ must satisfy a polynomial equation
\begin{equation*}
a _ {0} (z) w ^ {n} + a _ {1} (z) w ^ {n - 1} + \dots + a _ {n} (z) = 0,
\end{equation*}
and it is proved that $\mathbf{f}$ is algebraic.

It is now easy to settle the point which was left open in Sec.~\hyperref[sec:8.2.2]{2.2}. Suppose that the function element $(f,\Omega)$ satisfies the equation $P(f(z),z)=0$ where P is irreducible and of degree n in w. Then the corresponding global analytic\index{Function!analytic!global} function\index{Global analytic function} f has only algebraic singularities and a finite number of branches. According to what we have just shown f will satisfy a polynomial equation whose degree is equal to the number of branches. It will hence satisfy an irreducible equation whose degree is not higher. But the only irreducible equation it can satisfy is $P(w,z)=0$, and its degree is n. Therefore the number of branches is exactly n, and we have shown that all solutions of $P(w,z)=0$ are branches of the same analytic function.

It remains only to collect the results:

\begin{originalstatement}{Theorem}{4}\label{thm:8.4}
An analytic function is an algebraic function\index{Algebraic function}\index{Function!algebraic} if it has a finite number of branches and at most algebraic singularities. Every algebraic function $w = \mathbf{f}(z)$ satisfies an irreducible equation $P(w, z) = 0$, unique up to a constant factor, and every such equation determines a corresponding algebraic function uniquely.
\end{originalstatement}
It is also customary to say that an irreducible equation $P(w,z)=0$ defines an algebraic curve\index{Algebraic curve}. The theory of algebraic curves is a highly developed branch of algebra and function theory. We have been able to develop only the most elementary part of the function theoretic aspect.

\subsection*{EXERCISE}
\originalnumber{}{exercise:8.2.3.2.unnumbered}
Determine the position and nature of the singularities of the algebraic function defined by $w^{3} - 3wz + 2z^{3} = 0$.

\setcounter{section}{2}\section{Picard's Theorem}\index{Picard's theorem|(}\index{Picard's theorem}\label{sec:8.3}
In this section we shall prove the celebrated theorem of Picard, which asserts that an entire function omits at most one finite value. We shall
% Original PDF physical page 323
\index{Algebraic function|)}\index{Critical point|)}\index{Function!algebraic|)}\sourcepage{323}
prove it as an application of the monodromy theorem (Sec.~\hyperref[sec:8.1.6]{1.6}), using the modular function $\lambda(\tau)$ (Chap.~\hyperref[chap:7]{7}, Secs. 3.4 and 3.5) in an essential way. This is Picard's own proof. Many other proofs have been given which are more elementary in that they need less preparation, but none is as penetrating as the original proof.

\setcounter{section}{3}\setcounter{subsection}{0}\subsection{Lacunary Values}\label{sec:8.3.1}
A complex number $a$ is said to be a lacunary value\index{Lacunary value} of a function $f(z)$ if $f(z) \neq a$ in the region where $f$ is defined. For instance, 0 is a lacunary value of $e^z$ in the whole plane.

\begin{originalstatement}[Picard]{Theorem}{5}\label{thm:8.5}
An entire function with more than one finite lacunary value reduces to a constant.
\end{originalstatement}
We recall that an entire function $f(z)$ is one which is analytic in the whole plane. If a and b are distinct finite values and if $f(z)$ is different from a and b for all z, we are required to show that $f(z)$ is constant. Consider $f_{1}(z) = (f(z) - a)/(b - a)$. This function is entire and $\neq 0$ and 1. If $f_{1}$ is constant, so is f. Therefore it is no restriction to assume from the beginning that a = 0, b = 1.

We shall define a global analytic\index{Function!analytic!global} function\index{Global analytic function} $\mathbf{h}$ whose function elements $(h,\Omega)$ share the following property: $\operatorname{Im} h(z) > 0$, and $\lambda(h(z)) = f(z)$ for $z \in \Omega$. Here $\lambda(\tau)$ is the modular function defined in Chap.~\hyperref[chap:7]{7}, Sec.~\hyperref[sec:7.3.5]{3.5}. It will be shown that $\mathbf{h}$ can be continued along all paths. Since the plane is simply connected it will follow by the monodromy theorem that $\mathbf{h}$ defines an entire function $h(z)$. Because $h(z)$ has all its values in the upper half plane, $e^{ih}$ is bounded. By Liouville's theorem $h$ must reduce to a constant, and so does $f(z) = \lambda(h(z))$.

By Theorem~\hyperref[thm:7.7]{7} of Chap.~\hyperref[chap:7]{7} there exists a point $\tau_0$ in the upper half plane such that $\lambda(\tau_0) = f(0)$. Because $\lambda'(\tau_0) \neq 0$, by the same theorem, there exists a local inverse of $\lambda$, defined in a neighborhood $\Delta_0$ of $f(0)$ and denoted by $\lambda_0^{-1}$, characterized by the conditions $\lambda(\lambda_0^{-1}(w)) = w$ in $\Delta_0$ and
\begin{equation*}
\lambda_ {0} ^ {- 1} (f (0)) = \tau_ {0}.
\end{equation*}
By continuity there is a neighborhood $\Omega_0$ of the origin in which $f(z) \in \Delta_0$, and we can therefore define $h(z) = \lambda_0^{-1}(f(z))$ in $\Omega_0$. We shall let $\mathbf{h}$ be the global analytic function obtained by continuing the function element $(h, \Omega_0)$ in all possible ways.

We have to show that the element $(h,\Omega_{0})$ can be continued along all paths, and that Im h remains positive. If this were not so, we could find a path $\gamma[0,t_{1}]$ such that h can be continued and Im h remains positive up to any $t<t_{1}$, while either h cannot be continued up to $t_{1}$, or else Im $h[\gamma(t)]$ tends to 0 for $t\to t_{1}$. We can determine a value $\tau_{1}$ in the upper half-%
% Original PDF physical page 324
\sourcepage{324}%
plane with $\lambda(\tau_1) = f[\gamma(t_1)]$ and a local inverse $\lambda_1^{-1}$ with $\lambda_1^{-1}(f[\gamma(t_1)]) = \tau_1$, defined in a neighborhood $\Delta_1$ of $f[\gamma(t_1)]$. Let $\Omega_1$ be a neighborhood of $\gamma(t_1)$ in which $f(z) \in \Delta_1$, and choose $t_2 < t_1$ so that $\gamma(t) \in \Omega_1$ for $t \in [t_2, t_1]$. We know that $\lambda(\tau)$ has the same value $f[\gamma(t_2)]$ at $\tau = h[\gamma(t_2)]$ and at $\tau = \lambda_1^{-1}(f[\gamma(t_2)])$. Hence, by Theorem~\hyperref[thm:7.8]{8} of Chap.~\hyperref[chap:7]{7}, there exists a modular transformation $S$ in the congruence subgroup mod 2 such that
\begin{equation*}
S \left[ \lambda_ {1} ^ {- 1} (f [ \gamma (t _ {2}) ]) \right] = h [ \gamma (t _ {2}) ].
\end{equation*}
We now define $h_1$ in $\Omega_1$ by $h_1(z) = S[\lambda_1^{-1}(f(z))]$. It is evident that $(h_1, \Omega_1)$ is a continuation of $\mathbf{h}$ up to $t_1$ which satisfies $\lambda(h_1(z)) = f(z)$ and $\operatorname{Im} h_1 > 0$. We conclude that $\mathbf{h}$ can indeed be continued along all paths, and as we have pointed out, Picard's theorem\index{Picard's theorem} follows at once.

We have carried out the proof in such painstaking detail in an effort to convince the reader that the monodromy theorem plays as essential a role in the proof as the modular function.

\setcounter{section}{3}\section{Linear Differential Equations}\index{Linear differential equation}\label{sec:8.4}
The theory of global analytic\index{Function!analytic!global} functions\index{Global analytic function} makes it possible to study, with a great degree of generality, the complex solutions of ordinary differential equations\index{Differential equation|(}\index{Differential equation}. Of all differential equations the linear ones are the simplest, and also the most important. A linear equation of order n has the form
\begin{equation*}
a _ {0} (z) \frac {d ^ {n} w}{d z ^ {n}} + a _ {1} (z) \frac {d ^ {n - 1} w}{d z ^ {n - 1}} + \dots + a _ {n - 1} (z) \frac {d w}{d z} + a _ {n} (z) w = b (z)\tag{8}\label{eq:8.8}
\end{equation*}
where the coefficients $a_{k}(z)$ and the right-hand member $b(z)$ are single-valued analytic functions. In order to simplify the treatment we restrict our attention to the case where these functions are defined in the whole plane; they are thus assumed to be entire functions. A \emph{solution}\index{Solution} of \hyperref[eq:8.8]{(8)} is a global analytic function f which satisfies the identity
\begin{equation*}
a _ {0} \mathbf {f} ^ {(n)} + a _ {1} \mathbf {f} ^ {(n - 1)} + \dots + a _ {n - 1} \mathbf {f} ^ {\prime} + a _ {n} \mathbf {f} = b.\tag{9}\label{eq:8.9}
\end{equation*}
We have already remarked that this is a meaningful equation and that it is fulfilled as soon as a function element $(f,\Omega)$ of f satisfies the corresponding equation with f replaced by f. A function element with this property will be called a \emph{local solution}\index{Local solution}.

The reader who is familiar with the real case will expect the equation \hyperref[eq:8.9]{(9)} to have n linearly independent solutions. This is so as far as local solutions are concerned, but we must be prepared to find that different local solutions can be elements of the same global analytic function. In other words, in the complex case part of the problem is to find out to what extent the local solutions are analytic continuations of each other.

The equation \hyperref[eq:8.8]{(8)} is \emph{homogeneous} if $b(z)$ is identically zero. This is the most important case, and it is the only one we will treat. Further%
% Original PDF physical page 325
\index{Picard's theorem|)}\sourcepage{325}%
more, we can assume that the coefficients $a_{k}(z)$ have no common zeros; in fact, if $z_{0}$ were a common zero we could divide all coefficients by $z - z_{0}$, and the solutions would remain the same. As a matter of fact, if we are willing to consider meromorphic coefficients we may divide \hyperref[eq:8.8]{(8)} by $a_{0}(z)$ from the beginning. Conversely, if an equation with meromorphic coefficients is given, each coefficient can be written as a quotient of two entire functions; after multiplication with the common denominator we obtain an equivalent equation with entire coefficients. It is thus irrelevant whether we do or do not allow the coefficients to have poles.

In the case n = 1 the equation \hyperref[eq:8.8]{(8)} has the explicit solution
\begin{equation*}
w = e ^ {- \int \frac {a _ {1} (z)}{a _ {0} (z)} d z}.
\end{equation*}
The only problem is thus to determine the multiple-valued character of the integral, a question which has already been treated. On the other hand, the case n = 2 is found to have all the characteristic features of the general case. For this reason we find it sufficient to deal with homogeneous linear differential equations\index{Differential equation}\index{Linear differential equation} of the second order.

\setcounter{section}{4}\setcounter{subsection}{0}\subsection{Ordinary Points}\label{sec:8.4.1}
A point $z_0$ is called an \emph{ordinary point} for the differential equation
\begin{equation*}
a _ {0} (z) w ^ {\prime \prime} + a _ {1} (z) w ^ {\prime} + a _ {2} (z) w = 0\tag{10}\label{eq:8.10}
\end{equation*}
if and only if $a_0(z_0) \neq 0$. The central theorem to be proved is the following:

\begin{originalstatement}{Theorem}{6}\label{thm:8.6}
If $z_0$ is an ordinary point\index{Point!ordinary} for the equation\hyperref[eq:8.10]{(10)}, there exists a local solution $(f, \Omega)$, $z_0 \in \Omega$, with arbitrarily described values $f(z_0) = b_0$ and $f'(z_0) = b_1$. The germ $(f, z_0)$ is uniquely determined.
\end{originalstatement}
We prefer to write \hyperref[eq:8.10]{(10)} in the form
\begin{equation*}
w ^ {\prime \prime} = p (z) w ^ {\prime} + q (z) w\tag{11}\label{eq:8.11}
\end{equation*}
where $p(z) = -a_1 / a_0$, $q(z) = -a_2 / a_0$. The assumption means that $p(z)$ and $q(z)$ are analytic in a neighborhood of $z_0$; for convenience we may take $z_0 = 0$. Let
\begin{equation*}
\begin{array}{l} p (z) = p _ {0} + p _ {1} z + \dots + p _ {n} z ^ {n} + \dots \\ q (z) = q _ {0} + q _ {1} z + \dots + q _ {n} z ^ {n} + \dots \end{array}\tag{12}\label{eq:8.12}
\end{equation*}
be the Taylor developments of $p(z)$ and $q(z)$.

In order to solve \hyperref[eq:8.11]{(11)} we use the method of indeterminate coefficients. If the theorem is true, the solution $w = f(z)$ must have a Taylor development


% Original PDF physical page 326
\sourcepage{326}
\begin{equation*}
f (z) = b _ {0} + b _ {1} z + \dots + b _ {n} z ^ {n} + \dots\tag{13}\label{eq:8.13}
\end{equation*}
whose coefficients satisfy the conditions
\begin{equation*}
\begin{array}{l} 2 b _ {2} = b _ {1} p _ {0} + b _ {0} q _ {0} \\ 6 b _ {3} = 2 b _ {2} p _ {0} + b _ {1} p _ {1} + b _ {1} q _ {0} + b _ {0} q _ {1} \\ \dots \dots \\ n (n - 1) b _ {n} = (n - 1) b _ {n - 1} p _ {0} + (n - 2) b _ {n - 2} p _ {1} + \dots + b _ {1} p _ {n - 2} \\ \qquad \qquad \qquad + b _ {n - 2} q _ {0} + b _ {n - 3} q _ {1} + \dots + b _ {0} q _ {n - 2} \\ \dots \dots . \end{array}\tag{14}\label{eq:8.14}
\end{equation*}
This already proves the uniqueness. All that remains to prove is that the equations \hyperref[eq:8.14]{(14)} lead to a power series \hyperref[eq:8.13]{(13)} with a positive radius of convergence. It will then follow by permissible operations of term-wise differentiation, multiplication, and rearrangement that \hyperref[eq:8.13]{(13)} is a solution of the equation with desired initial values of f and $f'$.

Since the series \hyperref[eq:8.12]{(12)} have positive radii of convergence, there exist, by the Cauchy inequalities, constants $M_{0}$ and $r_{0} > 0$ such that
\begin{equation*}
\begin{array}{l l} | p _ {n} | \leq M _ {0} r _ {0} ^ {- n} \\ | q _ {n} | \leq M _ {0} r _ {0} ^ {- n}. \end{array}\tag{15}\label{eq:8.15}
\end{equation*}
In order to show that \hyperref[eq:8.13]{(13)} has likewise a positive radius of convergence, it is sufficient to prove similar inequalities
\begin{equation*}
\left| b _ {n} \right| \leq M r ^ {- n}\tag{16}\label{eq:8.16}
\end{equation*}
for a suitable choice of M and r.

The natural idea is to prove \hyperref[eq:8.16]{(16)} by induction on n. In the first place \hyperref[eq:8.16]{(16)} must hold for n = 0 and n = 1; this leads to the preliminary conditions $|b_{0}| \leq M$, $|b_{1}| \leq Mr^{-1}$ which are satisfied for sufficiently large M and sufficiently small r. Assume \hyperref[eq:8.16]{(16)} to be valid for all subscripts <n. In order to simplify the computations we choose $r < r_{0}$; then the general equation \hyperref[eq:8.14]{(14)} leads at once to the estimate
\begin{equation*}
\begin{array}{r l} n (n - 1) | b _ {n} | & \leq M M _ {0} [ (1 + 2 + \dots + (n - 1)) r ^ {1 - n} + (n - 1) r ^ {2 - n} ] \\ & = M M _ {0} \left[ \frac {n (n - 1)}{2} r + (n - 1) r ^ {2} \right] r ^ {- n}. \end{array}
\end{equation*}
We have thus
\begin{equation*}
\left| b _ {n} \right| \leq M M _ {0} \left(\frac {r}{2} + \frac {r ^ {2}}{n}\right) r ^ {- n} \leq M M _ {0} \left(\frac {r}{2} + r ^ {2}\right) r ^ {- n}
\end{equation*}
and \hyperref[eq:8.16]{(16)} follows, provided that $M_{0}(r/2 + r^{2}) \leq 1$. It is clear that this and the preceding requirements are fulfilled for all sufficiently small r. The proof is complete.

There exist, in particular, local solutions $f_0(z)$ and $f_1(z)$ which satisfy the conditions $f_0(z_0) = 1$, $f_0'(z_0) = 0$ and $f_1(z_0) = 0$, $f_1'(z_0) = 1$. Because
% Original PDF physical page 327
\sourcepage{327}
of the uniqueness the solution with the initial values $b_{0}$, $b_{1}$ must be $f(z) = b_{0}f_{0}(z) + b_{1}f_{1}(z)$. Hence every local solution is a linear combination of $f_{0}(z)$ and $f_{1}(z)$. Moreover, the solutions $f_{0}(z)$ and $f_{1}(z)$ are linearly independent, for if $b_{0}f_{0}(z) + b_{1}f_{1}(z) = 0$ we obtain first $b_{0} = 0$ by substituting $z = z_{0}$, and subsequently $b_{1} = 0$ since $f_{1}(z)$ cannot be identically zero.

\subsection*{EXERCISES}
\originalnumber{1}{exercise:8.4.1.3.1}
1. Find the power-series developments about the origin of two linearly independent solutions of $w'' = zw$.

\originalnumber{2}{exercise:8.4.1.3.2}
2. The Hermite polynomials are defined by
\begin{equation*}
H _ {n} (z) = (- 1) ^ {n} e ^ {z ^ {2}} \frac {d ^ {n}}{d z ^ {n}} \left(e ^ {- z ^ {2}}\right).
\end{equation*}
Prove that $H_{n}(z)$ is a solution of $w'' - 2zw' + 2nw = 0$.

\setcounter{section}{4}\setcounter{subsection}{1}\subsection{Regular Singular Points}\index{Regular singular point|(}\index{Regular singular point}\label{sec:8.4.2}
Any point $z_{0}$ such that $a_{0}(z_{0}) = 0$ is called a singular point\index{Singular point} of the equation \hyperref[eq:8.10]{(10)}. If the equation is written in the form \hyperref[eq:8.11]{(11)}, the assumption means that either $p(z)$ or $q(z)$ has a pole at $z_{0}$, for we continue to exclude the case of common zeros of all the coefficients in \hyperref[eq:8.10]{(10)}.

There are different kinds of singular points. We begin by a preliminary study of the simplest case which occurs when $a_{0}(z)$ has a simple zero. Under this hypothesis the functions $p(z)$ and $q(z)$ have at most simple poles, and if we choose $z_{0}=0$ the Laurent developments are of the form
\begin{equation*}
p (z) = \frac {p _ {- 1}}{z} + p _ {0} + p _ {1} z + \dots
\end{equation*}
\begin{equation*}
q (z) = \frac {q _ {- 1}}{z} + q _ {0} + q _ {1} z + \dots .
\end{equation*}
This time, if we substitute
\begin{equation*}
w = b _ {0} + b _ {1} z + b _ {2} z ^ {2} + \dots
\end{equation*}
in \hyperref[eq:8.11]{(11)}, the comparison of coefficients yields
\begin{equation*}
\begin{array}{l} - p _ {- 1} b _ {1} = b _ {0} q _ {- 1} \\ 2 (1 - p _ {- 1}) b _ {2} = b _ {1} p _ {0} + b _ {1} q _ {- 1} + b _ {0} q _ {0} \\ \dots \dots \\ n (n - 1 - p _ {- 1}) b _ {n} = (n - 1) b _ {n - 1} p _ {0} + (n - 2) b _ {n - 2} p _ {1} + \dots \\ \qquad + b _ {1} p _ {n - 2} + b _ {n - 1} q _ {- 1} + b _ {n - 2} q _ {0} + \dots + b _ {0} q _ {n - 2} \\ \dots \dots . \end{array}\tag{17}\label{eq:8.17}
\end{equation*}
% Original PDF physical page 328
\sourcepage{328}
This system of relations is essentially different from \hyperref[eq:8.14]{(14)}. In the first place, only $b_{0}$ can be chosen arbitrarily, and hence the method yields at most one linearly independent solution. Secondly, if $p_{-1}$ is zero or a positive integer, the system \hyperref[eq:8.17]{(17)} has either no solution or one of the $b_{n}$ can be chosen arbitrarily.

Assuming that $p_{-1}$ is not zero or a positive integer we will show that the resulting power series has a positive radius of convergence. As before we use the estimates \hyperref[eq:8.15]{(15)}, choose $M \geq |b_{0}|$, and assume \hyperref[eq:8.16]{(16)} for subscripts <n. Under the auxiliary hypothesis $r \leq r_{0}$ we obtain
\begin{equation*}
\begin{aligned}
&n|n-1-p_{-1}|\cdot|b_n|\\
&\quad\leq Mr^{-n}\left\{M_0\left[\frac{n(n-1)}2r+(n-1)r^2\right]+|q_{-1}|r\right\}.
\end{aligned}
\end{equation*}
Inasmuch as $(n - 1) / |n - 1 - p_{-1}|$ is bounded, an inequality of the form
\begin{equation*}
\left| b _ {n} \right| \leq M r ^ {- n} (A r + B r ^ {2})
\end{equation*}
will hold for all n. For sufficiently small r this is stronger than \hyperref[eq:8.16]{(16)}, and the convergence follows.

As already indicated, the result is of a preliminary nature. Our real object is to solve \hyperref[eq:8.11]{(11)} in the presence of a \emph{regular singularity} at $z_{0}$. This terminology is used to indicate that $p(z)$ has at most a simple and $q(z)$ at most a double pole at $z_{0}$.

Under these circumstances it turns out that there are solutions of the form $w = z^{\alpha}g(z)$ where $g(z)$ is analytic and $\neq0$ at $z_{0}(=0)$. We make this substitution in \hyperref[eq:8.11]{(11)} and find, after brief computation, that $g(z)$ must satisfy the differential equation\index{Differential equation}
\begin{equation*}
g ^ {\prime \prime} = \left(p - \frac {2 \alpha}{z}\right) g ^ {\prime} + \left(q + \frac {\alpha p}{z} - \frac {\alpha (\alpha - 1)}{z ^ {2}}\right) g.\tag{18}\label{eq:8.18}
\end{equation*}
For arbitrary $\alpha$ this is of the same type as the original equation, and nothing has been gained. We may, however, choose $\alpha$ so that the coefficient of g has only a simple pole. If $q(z)$ has the development
\begin{equation*}
q (z) = \frac {q _ {- 2}}{z ^ {2}} + \dots
\end{equation*}
this will be the case if $\alpha$ satisfies the quadratic equation
\begin{equation*}
\alpha (\alpha - 1) - p _ {- 1} \alpha - q _ {- 2} = 0,\tag{19}\label{eq:8.19}
\end{equation*}
known as the \emph{indicial equation}\index{Indicial equation}. For such $\alpha$ our preliminary result shows that \hyperref[eq:8.11]{(11)} has a solution of the form $z^{\alpha}g(z)$, $g(0) \neq 0$, provided that $p_{-1} - 2\alpha$ is not a nonnegative integer.

Let the roots of \hyperref[eq:8.19]{(19)} be denoted by $\alpha_{1}$ and $\alpha_{2}$. Then


% Original PDF physical page 329
\sourcepage{329}
\begin{equation*}
\alpha_ {1} + \alpha_ {2} = p _ {- 1} + 1
\end{equation*}
or $\alpha_{2}-\alpha_{1}=p_{-1}-2\alpha_{1}+1$. Hence $\alpha_{1}$ is exceptional if and only if $\alpha_{2}-\alpha_{1}$ is a positive integer; by symmetry, $\alpha_{2}$ is exceptional if $\alpha_{2}-\alpha_{1}$ is a negative integer. Consequently, if the roots of the indicial equation do not differ by an integer, we obtain two solutions $z^{\alpha_{1}}g_{1}(z)$ and $z^{\alpha_{2}}g_{2}(z)$ which are obviously linearly independent. If the roots are equal or differ by an integer, the method yields only one solution.

\begin{originalstatement}{Theorem}{7}\label{thm:8.7}
If $z_0$ is a regular singular point\index{Regular singular point} for the equation \hyperref[eq:8.10]{(10)}, there exist linearly independent solutions of the form $(z - z_0)^{\alpha_1}g_1(z)$ and $(z - z_0)^{\alpha_2}g_2(z)$ with $g_1(0), g_2(0) \neq 0$ corresponding to the roots of the indicial equation, provided that $\alpha_2 - \alpha_1$ is not an integer. In the case of an integral difference $\alpha_2 - \alpha_1 \geq 0$ the existence of a solution corresponding to $\alpha_2$ can still be asserted.
\end{originalstatement}
If one solution is known it is not difficult to find another, linearly independent of the first. The methods which lead to a second solution belong more properly in a textbook on differential equations\index{Differential equation}. It is also impossible to treat the case of irregular singularities in this book.

\subsection*{EXERCISES}
\originalnumber{1}{exercise:8.4.2.4.1}
1. Show that the equation $(1 - z^{2})w'' - 2zw' + n(n + 1)w = 0$, where $n$ is a nonnegative integer, has the Legendre polynomials
\begin{equation*}
P _ {n} (z) = \frac {1}{2 ^ {n} n !} \cdot \frac {d ^ {n}}{d z ^ {n}} (z ^ {2} - 1) ^ {n}
\end{equation*}
as solutions.

\originalnumber{2}{exercise:8.4.2.4.2}
2. Determine two linearly independent solutions of the equation
\begin{equation*}
z ^ {2} (z + 1) w ^ {\prime \prime} - z ^ {2} w ^ {\prime} + w = 0
\end{equation*}
near 0 and one near -1.

\originalnumber{3}{exercise:8.4.2.4.3}
3. Show that Bessel\index{Bessel F W@Bessel, F. W.}'s equation $zw'' + w' + zw = 0$ has a solution which is an integral function. Determine its power-series development.

\setcounter{section}{4}\setcounter{subsection}{2}\subsection{Solutions at Infinity}\label{sec:8.4.3}
If $a_{0}(z)$, $a_{1}(z)$, $a_{2}(z)$ are polynomials, it is natural to ask how the solutions behave in the neighborhood of $\infty$. The most convenient way to treat this question is to make the variable transformation z = 1/Z. Since
\begin{equation*}
\frac {d w}{d z} = - Z ^ {2} \frac {d w}{d Z}
\end{equation*}
\begin{equation*}
\frac {d ^ {2} w}{d z ^ {2}} = 2 Z ^ {3} \frac {d w}{d Z} + Z ^ {4} \frac {d ^ {2} w}{d Z ^ {2}}
\end{equation*}
% Original PDF physical page 330
\index{Regular singular point|)}\sourcepage{330}
equation \hyperref[eq:8.11]{(11)} takes the form
\begin{equation*}
\frac {d ^ {2} w}{d Z ^ {2}} = - \left(2 Z ^ {- 1} + Z ^ {- 2} p \left(\frac {1}{Z}\right)\right) \frac {d w}{d Z} + Z ^ {- 4} q \left(\frac {1}{Z}\right) w.\tag{20}\label{eq:8.20}
\end{equation*}
We say of course that $\infty$ is an ordinary point or a regular singularity for the equation \hyperref[eq:8.11]{(11)} if the point Z = 0 has the corresponding character for \hyperref[eq:8.20]{(20)}. Thus $\infty$ is an ordinary point if the coefficients in \hyperref[eq:8.11]{(11)} have a removable singularity at Z = 0; this is the same, by definition, as saying that $-(2z + z^{2}p(z))$ and $z^{4}q(z)$ have removable singularities at $\infty$. Similarly, $\infty$ is a regular singularity if these functions have, respectively, at most a simple and a double pole at $\infty$.

It is interesting to determine the equations with the fewest singularities. If $\infty$ is to be an ordinary point, $q(z)$ must have at least four poles, unless it vanishes identically. In the latter case $p(z)$ can have as few as one pole, and if the pole is placed at the origin we must have $p(z) = -2/z$. The corresponding equation
\begin{equation*}
\frac {d ^ {2} w}{d z ^ {2}} = - \frac {2}{z} \frac {d w}{d z}
\end{equation*}
has the general solution $w = az^{-1} + b$.

If $q(z)$ is not identically zero, there can be as few as two regular singularities. It is evidently easiest to place the singularities at 0 and $\infty$, and for this reason we turn immediately to the case where $\infty$ is a regular singularity. If there is to be only one finite singularity, placed at the origin, we must have $p(z) = A/z$, $q(z) = B/z^{2}$. With another choice of constants the equation can be written in the form
\begin{equation*}
z ^ {2} w ^ {\prime \prime} - (\alpha + \beta - 1) z w ^ {\prime} + \alpha \beta w = 0.\tag{21}\label{eq:8.21}
\end{equation*}
It has the solutions $w = z^{\alpha}$ and $w = z^{\beta}$, where $\alpha$ and $\beta$ are obviously the roots of the indicial equation. If $\alpha = \beta$, there must be another solution. To find it we write \hyperref[eq:8.21]{(21)} in the symbolic form
\begin{equation*}
\left(z \frac {d}{d z} - \alpha\right) ^ {2} w = 0
\end{equation*}
and substitute $w = z^{\alpha}W$. We obtain
\begin{equation*}
\left(z \frac {d}{d z} - \alpha\right) z ^ {\alpha} W = z ^ {\alpha} \cdot z \frac {d W}{d z}
\end{equation*}
\begin{equation*}
\left(z \frac {d}{d z} - \alpha\right) ^ {2} z ^ {\alpha} W = z ^ {\alpha} \cdot z \frac {d}{d z} \left(z \frac {d W}{d z}\right).
\end{equation*}
% Original PDF physical page 331
\sourcepage{331}
The\index{Function!hypergeometric|(} equation $\left(z\frac{d}{dz}\right)^2 W = 0$ has the obvious solution $W = \log z$, and hence the desired solution of \hyperref[eq:8.21]{(21)} is $w = z^{\alpha} \log z$.

\setcounter{section}{4}\setcounter{subsection}{3}\subsection{The Hypergeometric Differential Equation}\index{Hypergeometric differential equation|(}\index{Hypergeometric differential equation}\label{sec:8.4.4}
We have just seen that differential equations\index{Differential equation} with one or two regular singularities have trivial solutions. It is only with the introduction of a third singularity that we obtain a new and interesting class of analytic functions.

It is quite clear that a linear transformation of the variable transforms a second-order linear differential equation\index{Linear differential equation} into one of the same type and that the character of the singularities remains the same. We can therefore elect to place the three singularities at prescribed points, and it is simplest to choose them at 0, 1, and $\infty$.

If the equation
\begin{equation*}
w ^ {\prime \prime} = p (z) w ^ {\prime} + q (z)
\end{equation*}
is to have finite regular singularities only at 0 and 1, we must have
\begin{equation*}
\begin{array}{l} p (z) = \frac {A}{z} + \frac {B}{z - 1} + P (z) \\ q (z) = \frac {C}{z ^ {2}} + \frac {D}{z} + \frac {E}{(z - 1) ^ {2}} + \frac {F}{z - 1} + Q (z) \end{array}
\end{equation*}
where $P(z)$ and $Q(z)$ are polynomials. In order to make the singularity at $\infty$ regular, $2z + z^{2}p(z)$ must have at most a simple pole at $\infty$ and $z^{4}q(z)$ must have at most a double pole. In view of these conditions $P(z)$ and $Q(z)$ must be identically zero, and the relation $D + F = 0$ must hold. These are evidently the only conditions, and we can rewrite the expressions for $p(z)$ and $q(z)$ in the form
\begin{equation*}
\begin{array}{l} p (z) = \frac {A}{z} + \frac {B}{z - 1} \\ q (z) = \frac {C}{z ^ {2}} - \frac {D}{z (z - 1)} + \frac {E}{(z - 1) ^ {2}}. \end{array}
\end{equation*}
The indicial equation at the origin reads
\begin{equation*}
\alpha (\alpha - 1) = A \alpha + C.
\end{equation*}
So if its roots are denoted by $\alpha_{1}, \alpha_{2}$ we obtain $A = \alpha_{1} + \alpha_{2} - 1$, $C = -\alpha_{1}\alpha_{2}$. Similarly, $B = \beta_{1} + \beta_{2} - 1$ and $E = -\beta_{1}\beta_{2}$, where $\beta_{1}, \beta_{2}$ are the roots of the indicial equation at 1. In order to write down the indicial equation at $\infty$ we note that the leading coefficients of $-2z -$
% Original PDF physical page 332
\sourcepage{332}
$z^{2}p(z)$ and $z^{4}q(z)$ are $-(2+A+B)$ and $C-D+E$, respectively. Hence the roots $\gamma_{1},\gamma_{2}$ satisfy $\gamma_{1}+\gamma_{2}=-A-B-1$ and
\begin{equation*}
\gamma_ {1} \gamma_ {2} = - C + D - E.
\end{equation*}
We conclude at the relation
\begin{equation*}
\alpha_ {1} + \alpha_ {2} + \beta_ {1} + \beta_ {2} + \gamma_ {1} + \gamma_ {2} = 1,\tag{22}\label{eq:8.22}
\end{equation*}
and we find that the equation can be written in the form
\begin{equation*}
\begin{aligned}
&w''+\left(\frac{1-\alpha_1-\alpha_2}{z}+\frac{1-\beta_1-\beta_2}{z-1}\right)w'\\
&\quad+\left(\frac{\alpha_1\alpha_2}{z^2}
-\frac{\alpha_1\alpha_2+\beta_1\beta_2-\gamma_1\gamma_2}{z(z-1)}
+\frac{\beta_1\beta_2}{(z-1)^2}\right)w=0.
\end{aligned}\tag{23}\label{eq:8.23}
\end{equation*}
In order to avoid the exceptional cases we will now assume that none of the differences $\alpha_{2}-\alpha_{1},\beta_{2}-\beta_{1},\gamma_{2}-\gamma_{1}$ is an integer. Our next step is to simplify the equation \hyperref[eq:8.23]{(23)}. In Sec.~\hyperref[sec:8.4.2]{4.2} we have already shown that the substitution $w=z^{\alpha}g(z)$ determines for $g(z)$ a similar differential equation\index{Differential equation}, namely, the equation \hyperref[eq:8.18]{(18)}. Since the original equation has solutions of the form $w=z^{\alpha_{1}}g_{1}(z),w=z^{\alpha_{2}}g_{2}(z)$, we conclude that the transformed equation \hyperref[eq:8.18]{(18)} must have solutions of the form $g(z)=z^{\alpha_{1}-\alpha}g_{1}(z)$ and $g(z)=z^{\alpha_{2}-\alpha}g_{2}(z)$. Hence the indicial equation of \hyperref[eq:8.18]{(18)} has the roots $\alpha_{1}-\alpha,\alpha_{2}-\alpha$, as can also be verified by computation. Simultaneously, the roots which correspond to the singularity at $\infty$ change from $\gamma_{1},\gamma_{2}$ to $\gamma_{1}+\alpha,\gamma_{2}+\alpha$. In exactly the same way we can separate a factor $(z-1)^{\beta}$ and find that the resulting equation has exponents which are smaller by $\beta$ at 1 and larger by $\beta$ at $\infty$. The natural choice is to take $\alpha=\alpha_{1},\beta=\beta_{1}$. In the final equation the six exponents are then 0, $\alpha_{2}-\alpha_{1},0,\beta_{2}-\beta_{1},\gamma_{1}+\alpha_{1}+\beta_{1},\gamma_{2}+\alpha_{1}+\beta_{1}$, respectively. In order to comply with time-honored conventions we will write $a=\alpha_{1}+\beta_{1}+\gamma_{1},b=\alpha_{1}+\beta_{1}+\gamma_{2},c=1+\alpha_{1}-\alpha_{2}$. Because of the relation \hyperref[eq:8.22]{(22)} we get c-a-b= $\beta_{2}-\beta_{1}$. Accordingly, the new differential equation will be of the form
\begin{equation*}
w ^ {\prime \prime} + \left(\frac {c}{z} + \frac {1 - c + a + b}{z - 1}\right) w ^ {\prime} + \frac {a b}{z (z - 1)} w = 0
\end{equation*}
or, after simplification,
\begin{equation*}
z (1 - z) w ^ {\prime \prime} + [ c - (a + b + 1) z ] w ^ {\prime} - a b w = 0.\tag{24}\label{eq:8.24}
\end{equation*}
This is called the \emph{hypergeometric differential equation}\index{Hypergeometric differential equation}, and we have proved that the solutions of \hyperref[eq:8.23]{(23)} are equal to the solutions of \hyperref[eq:8.24]{(24)} multiplied by $z^{\alpha_{1}}(z - 1)^{\beta_{1}}$. It is assumed that none of the exponent differences c - 1,


% Original PDF physical page 333
\sourcepage{333}
$a - b, a + b - c$ is an integer.

According to the theory, equation \hyperref[eq:8.24]{(24)} has a solution of the form $w = \sum_{n=0}^{\infty} A_{n} z^{n}$. If this power series is substituted in \hyperref[eq:8.24]{(24)}, we find with very little computation that the coefficients must satisfy the recursive relations
\begin{equation*}
(n + 1) (n + c) A _ {n + 1} = (n + a) (n + b) A _ {n}.
\end{equation*}
The extremely simple form of this relation makes it possible to write down the solution explicitly. With the choice $A_{0} = 1$ we find that the hypergeometric equation is satisfied by the function
\begin{equation*}
\begin{aligned}
F(a,b,c,z)={}&1+\frac{a\cdot b}{1\cdot c}z
+\frac{a(a+1)\cdot b(b+1)}{1\cdot2\cdot c(c+1)}z^2\\
&+\frac{a(a+1)(a+2)\cdot b(b+1)(b+2)}{1\cdot2\cdot3\cdot c(c+1)(c+2)}z^3+\dots,
\end{aligned}
\end{equation*}
known as the \emph{hypergeometric function}\index{Function!hypergeometric}\index{Hypergeometric function}. It is defined as soon as c is not zero or a negative integer.

The radius of convergence of the hypergeometric series can easily be found by computation, but it is more instructive to use pure reasoning. In the first place, we know that $F(a,b,c,z)$ can be continued analytically along any path which does not pass through the point 1 and does not return to the origin. Hence a single-valued branch of $F(a,b,c,z)$ can be defined in the unit disk $|z| < 1$ (because the disk is simply connected), and it follows that the radius of convergence is at least equal to one. If it is greater than one, $F(a,b,c,z)$ will be an entire function. Near infinity it must be a linear combination of the solutions $z^{-a}g_1(z)$, $z^{-b}g_2(z)$ known to exist in a neighborhood of $\infty$. But it is clear that a linear combination can be single-valued only if $a$ or $b$ is an integer. If $a$ is an integer $b$ is not, by assumption, and $F(a,b,c,z)$ is a multiple of $z^{-a}g_1(z)$. By Liouville's theorem, if $a$ were positive $F(a,b,c,z)$ would vanish identically, which is not the case. The only case in which the radius of convergence is infinite is thus when $a$ (or $b$) is a negative integer or zero, and then the hypergeometric series reduces trivially to a polynomial.

In a neighborhood of the origin there is also a solution of the form $z^{1-c}g(z)$. Here $g(z)$ satisfies a hypergeometric differential equation\index{Differential equation}\index{Hypergeometric differential equation} with the six exponents $\alpha_{2}-\alpha_{1},0,0,\beta_{2}-\beta_{1},\gamma_{1}+\alpha_{2}+\beta_{1},\gamma_{2}+\alpha_{2}+\beta_{1}$. It follows at once that we can set $g(z)=F(1+a-c,1+b-c,2-c,z)$. We have proved that two linearly independent solutions near the origin are $F(a,b,c,z)$ and $z^{1-c}F(1+a-c,1+b-c,2-c,z)$, respectively.

The solutions near 1 can be determined in exactly the same manner. It is easier, however, to replace $z$ by $1 - z$ and interchange the $\alpha$'s and $\beta$'s.


% Original PDF physical page 334
\sourcepage{334}
As a result we find that the functions $F(a,b,1+a+b-c,1-z)$ and $(1-z)^{c-a-b}F(c-b,c-a,1-a-b+c,1-z)$ are linearly independent solutions in a neighborhood of 1. The solutions near $\infty$ can be found similarly.

We have demonstrated that the most general linear second-order differential equation\index{Differential equation} with three regular singularities can be solved explicitly by means of the hypergeometric function\index{Function!hypergeometric}. It is evidently also possible, although somewhat laborious, to determine the complete multiple-valued structure of the solutions.

\subsection*{EXERCISES}
\originalnumber{1}{exercise:8.4.4.5.1}
1. Show that $(1 - z)^{-\alpha} = F(\alpha, \beta, \beta, z)$ and $\log 1 / (1 - z) = zF(1, 1, 2, z)$.

\originalnumber{2}{exercise:8.4.4.5.2}
2. Express the derivative of $F(a,b,c,z)$ as a hypergeometric function.

\originalnumber{3}{exercise:8.4.4.5.3}
3. Derive the integral representation
\begin{equation*}
F (a, b, c, z) = \frac {\Gamma (c)}{\Gamma (b) \Gamma (c - b)} \int_ {0} ^ {1} t ^ {b - 1} (1 - t) ^ {c - b - 1} (1 - z t) ^ {- a} d t.
\end{equation*}
\originalnumber{4}{exercise:8.4.4.5.4}
4. If $w_{1}$ and $w_{2}$ are linearly independent solutions of the differential equation $w'' = pw' + qw$, prove that the quotient $\eta = w_{2}/w_{1}$ satisfies
\begin{equation*}
\frac {d}{d z} \left(\frac {\eta^ {\prime \prime}}{\eta^ {\prime}}\right) - \frac {1}{2} \left(\frac {\eta^ {\prime \prime}}{\eta^ {\prime}}\right) ^ {2} = - 2 q - \frac {1}{2} p ^ {2} + p ^ {\prime}.
\end{equation*}
\setcounter{section}{4}\setcounter{subsection}{4}\subsection{Riemann's Point of View}\label{sec:8.4.5}
Riemann was a strong proponent of the idea that an analytic function can be defined by its singularities and general properties just as well as or perhaps better than through an explicit expression. A trivial example is the determination of a rational function by the singular parts connected with its poles.

We will show, with Riemann, that the solutions of a hypergeometric differential equation\index{Hypergeometric differential equation} can be characterized by properties of this nature. We consider in the following a collection F of function elements $(f,\Omega)$ with certain characteristic features which we proceed to enumerate.

1. The collection $\mathbf{F}$ is \emph{complete} in the sense that it contains all analytic continuations of any $(f,\Omega)\in\mathbf{F}$. It is not required that any two function elements in $\mathbf{F}$ be analytic continuations of each other, and hence $\mathbf{F}$ may consist of several global analytic\index{Function!analytic!global} functions\index{Global analytic function}.

2. The collection is \emph{linear}. This means that $(f_{1},\Omega)\in\mathbf{F}$, $(f_{2},\Omega)\in\mathbf{F}$ implies $(c_{1}f_{1}+c_{2}f_{2},\Omega)\in\mathbf{F}$ for all constant $c_{1}$, $c_{2}$. Moreover, any three elements $(f_{1},\Omega)$, $(f_{2},\Omega)$, $(f_{3},\Omega)\in\mathbf{F}$ with the same $\Omega$ shall satisfy an identical relation $c_{1}f_{1}+c_{2}f_{2}+c_{3}f_{3}=0$ in $\Omega$ with constant coefficients, not all zero. In other words, F shall be at most \emph{two dimensional}.

3. The only finite singularities of the functions in F shall be at the
% Original PDF physical page 335
\sourcepage{335}
points 0 and 1; in addition, the point $\infty$ is also counted as a singularity. More precisely, it is required that any $(f,\Omega)\in\mathbf{F}$ can be continued along all arcs in the finite plane which do not pass through the points 0 and 1.

4. As to the behavior at the singular points we assume that there are functions in $\mathbf{F}$ which behave like prescribed powers $z^{\alpha_1}$ and $z^{\alpha_2}$ near 0, like $(z - 1)^{\beta_1}$ and $(z - 1)^{\beta_2}$ near 1, and like $z^{-\gamma_1}$ and $z^{-\gamma_2}$ near $\infty$. In precise terms, there shall exist certain analytic functions $g_1(z)$ and $g_2(z)$ defined in a neighborhood $\Delta$ of 0 and different from zero at that point; for a simply connected subregion $\Omega$ of $\Delta$ which does not contain the origin function elements $(z^{\alpha_1}g_1(z),\Omega)$, $(z^{\alpha_2}g_2(z),\Omega)$ can be defined, and it is required that they belong to $\mathbf{F}$. The corresponding assumptions for the points 1 and $\infty$ can be formulated in analogous manner.

The reader will have recognized that the solutions of the differential equation\index{Differential equation} \hyperref[eq:8.23]{(23)} have just these properties, provided that none of the differences $\alpha_{2}-\alpha_{1},\beta_{2}-\beta_{1},\gamma_{2}-\gamma_{1}$ is an integer. In addition, the relation $\alpha_{1}+\alpha_{2}+\beta_{1}+\beta_{2}+\gamma_{1}+\gamma_{2}=1$ is satisfied. We make both assumptions and prove, under these restrictions, that there exists one and only one collection F with the properties 1 to 4. Accordingly, F will be identical with the collection of local solutions of the differential equation \hyperref[eq:8.23]{(23)}.

Riemann denotes any function element in $\mathbf{F}$ by the symbol
\begin{equation*}
P \left\{ \begin{array}{c c c} 0 & 1 & \infty \\ \alpha_ {1} & \beta_ {1} & \gamma_ {1}, z \\ \alpha_ {2} & \beta_ {2} & \gamma_ {2} \end{array} \right\}.
\end{equation*}
Thus P does not stand for an individual function, but this is evidently of little importance. Once the uniqueness is established such identities as
\begin{equation*}
\begin{aligned}
&P\left\{\begin{array}{ccc}
0&1&\infty\\\alpha_1&\beta_1&\gamma_1,z\\\alpha_2&\beta_2&\gamma_2
\end{array}\right\}\\
&\quad=z^\alpha(z-1)^\beta P\left\{\begin{array}{ccc}
0&1&\infty\\\alpha_1-\alpha&\beta_1-\beta&\gamma_1+\alpha+\beta,z\\
\alpha_2-\alpha&\beta_2-\beta&\gamma_2+\alpha+\beta
\end{array}\right\}
\end{aligned}
\end{equation*}
or
\begin{equation*}
P \left\{ \begin{array}{c c c} 0 & 1 & \infty \\ \alpha_ {1} & \beta_ {1} & \gamma_ {1}, z \\ \alpha_ {2} & \beta_ {2} & \gamma_ {2} \end{array} \right\} = P \left\{ \begin{array}{c c c} 0 & 1 & \infty \\ \beta_ {1} & \alpha_ {1} & \gamma_ {1}, 1 - z \\ \beta_ {2} & \alpha_ {2} & \gamma_ {2} \end{array} \right\}
\end{equation*}
follow immediately provided that some care is given to their proper interpretation. The fact that such relationships, some of them quite elaborate, can be so easily recognized is one of the motivations for Riemann's point of view.

In order to prove the uniqueness, consider two linearly independent function elements $(f_{1},\Omega)$, $(f_{2},\Omega)\in\mathbf{F}$, defined in a simply connected region $\Omega$ which does not contain 0 or 1. There are such function elements in
% Original PDF physical page 336
\sourcepage{336}
any $\Omega$, for the functions $z^{\alpha_1}g_1(z)$ and $z^{\alpha_2}g_2(z)$ are linearly independent in their common region of definition; they can be continued along an arc that avoids 0 and 1 and ends in $\Omega$, where the continuations define linearly independent function elements $(f_1,\Omega)$, $(f_2,\Omega)$. By property 1 they belong to $\mathbf{F}$. If $(f,\Omega)$ is a third function element in $\mathbf{F}$, the identities
\begin{equation*}
\begin{array}{r l} {c f} & {+ c _ {1} f _ {1} + c _ {2} f _ {2} = 0} \\ {c f ^ {\prime}} & {+ c _ {1} f _ {1} ^ {\prime} + c _ {2} f _ {2} ^ {\prime} = 0} \\ {c f ^ {\prime \prime}} & {+ c _ {1} f _ {1} ^ {\prime \prime} + c _ {2} f _ {2} ^ {\prime \prime} = 0} \end{array}
\end{equation*}
imply
\begin{equation*}
\left| \begin{array}{c c c} f & f _ {1} & f _ {2} \\ f ^ {\prime} & f _ {1} ^ {\prime} & f _ {2} ^ {\prime} \\ f ^ {\prime \prime} & f _ {1} ^ {\prime \prime} & f _ {2} ^ {\prime \prime} \end{array} \right| = 0.
\end{equation*}
We write this equation in the form
\begin{equation*}
f ^ {\prime \prime} = p (z) f ^ {\prime} + q (z) f
\end{equation*}
with
\begin{equation*}
p (z) = \frac {f _ {1} f _ {2} ^ {\prime \prime} - f _ {2} f _ {1} ^ {\prime \prime}}{f _ {1} f _ {2} ^ {\prime} - f _ {2} f _ {1} ^ {\prime}}, \quad q (z) = - \frac {f _ {1} ^ {\prime} f _ {2} ^ {\prime \prime} - f _ {2} ^ {\prime} f _ {1} ^ {\prime \prime}}{f _ {1} f _ {2} ^ {\prime} - f _ {2} f _ {1} ^ {\prime}}.\tag{25}\label{eq:8.25}
\end{equation*}
Here the denominator is not identically zero, for that would mean that $f_{1}$ and $f_{2}$ were linearly dependent.

We make now the observation that the expressions \hyperref[eq:8.25]{(25)} remain invariant if $f_{1}$ and $f_{2}$ are subjected to a nonsingular linear transformation, i.e., if they are replaced by $c_{11}f_{1} + c_{12}f_{2}$, $c_{21}f_{1} + c_{22}f_{2}$ with $c_{11}c_{22} - c_{12}c_{21} \neq 0$. This means that $p(z)$ and $q(z)$ will be the same for any choice of $f_{1}$ and $f_{2}$; hence they are well-determined single-valued functions in the whole plane minus the points 0 and 1.

In order to determine the behavior of $p(z)$ and $q(z)$ near the origin, we choose $f_{1} = z^{\alpha_{1}}g_{1}(z), f_{2} = z^{\alpha_{2}}g_{2}(z)$. Simple calculations give
\begin{equation*}
\begin{array}{r l} f _ {1} f _ {2} ^ {\prime} - f _ {2} f _ {1} ^ {\prime} & = (\alpha_ {2} - \alpha_ {1}) z ^ {\alpha_ {1} + \alpha_ {2} - 1} (C + \dots) \\ f _ {1} f _ {2} ^ {\prime \prime} - f _ {2} f _ {1} ^ {\prime \prime} & = (\alpha_ {2} - \alpha_ {1}) (\alpha_ {1} + \alpha_ {2} - 1) z ^ {\alpha_ {1} + \alpha_ {2} - 2} (C + \dots) \\ f _ {1} ^ {\prime} f _ {2} ^ {\prime \prime} - f _ {2} ^ {\prime} f _ {1} ^ {\prime \prime} & = \alpha_ {1} \alpha_ {2} (\alpha_ {2} - \alpha_ {1}) z ^ {\alpha_ {1} + \alpha_ {2} - 3} (C + \dots) \end{array}
\end{equation*}
where the parentheses stand for analytic functions with the common value $C = g_{1}(0)g_{2}(0)$ at the origin. We conclude that $p(z)$ has a simple pole with the residue $\alpha_{1} + \alpha_{2} - 1$ while the Laurent development of $q(z)$ begins with the term $-\alpha_{1}\alpha_{2}/z^{2}$. Similar results hold for the points 1 and $\infty$. We infer that


% Original PDF physical page 337
\sourcepage{337}
\begin{equation*}
p (z) = \frac {\alpha_ {1} + \alpha_ {2} - 1}{z} + \frac {\beta_ {1} + \beta_ {2} - 1}{z - 1} + p _ {0} (z)\tag{26}\label{eq:8.26}
\end{equation*}
where $p_{0}(z)$ is free from poles at 0 and 1. According to its definition \hyperref[eq:8.24]{(24)}, $p(z)$ is the logarithmic derivative of an entire function; as such it has, in the finite plane, only simple poles with positive integers as residues. Moreover, the development of $p(z)$ at $\infty$ must begin with the term $-(\gamma_{1} + \gamma_{2} + 1)/z$. Hence $p(z)$ has only finitely many poles, and their residues must add up to $-(\gamma_{1} + \gamma_{2} + 1)$. In view of the relation $(\alpha_{1} + \alpha_{2} - 1) + (\beta_{1} + \beta_{2} - 1) = -(\gamma_{1} + \gamma_{2} + 1)$, it follows that there are no poles other than the ones at 0 and 1. A look at \hyperref[eq:8.26]{(26)} shows that $p_{0}(z)$ is pole-free and zero at $\infty$, hence identically zero.

Since $f_{1}f_{2}^{\prime} - f_{2}f_{1}^{\prime} \neq 0$ except at 0 and 1, we conclude that $q(z)$ is of the form
\begin{equation*}
q (z) = - \frac {\alpha_ {1} \alpha_ {2}}{z ^ {2}} - \frac {\beta_ {1} \beta_ {2}}{(z - 1) ^ {2}} + \frac {A}{z} + \frac {B}{z - 1} + q _ {0} (z)
\end{equation*}
where $q_0(z)$ has no finite poles. At $\infty$ the development must begin with $-\gamma_1\gamma_2 / z^2$. It follows that
\begin{equation*}
A = - B = - \left(\alpha_ {1} \alpha_ {2} + \beta_ {1} \beta_ {2} - \gamma_ {1} \gamma_ {2}\right)
\end{equation*}
and that $q_0(z)$ is identically zero. Collecting the results we conclude that $f$ is a solution of the equation
\begin{equation*}
\begin{aligned}
&w''+\left(\frac{1-\alpha_1-\alpha_2}{z}+\frac{1-\beta_1-\beta_2}{z-1}\right)w'\\
&\quad+\left(\frac{\alpha_1\alpha_2}{z^2}
-\frac{\alpha_1\alpha_2+\beta_1\beta_2-\gamma_1\gamma_2}{z(z-1)}
+\frac{\gamma_1\gamma_2}{(z-1)^2}\right)w=0
\end{aligned}
\end{equation*}
which is just equation \hyperref[eq:8.23]{(23)}.

This completes the uniqueness proofs, for it follows now that any collection $\mathbf{F}$ which satisfies 1 to 4 must be a subcollection of the family $\mathbf{F}_0$ of local solutions of \hyperref[eq:8.23]{(23)}. For any simply connected $\Omega$ which does not contain 0 or 1 we know that there are two linearly independent function elements $(f_1,\Omega)$, $(f_2,\Omega)$ in $\mathbf{F}$. Every $(f,\Omega) \in \mathbf{F}_0$ is of the form $(c_1f_1 + c_2f_2,\Omega)$ and is consequently contained in $\mathbf{F}$. Finally, if $\Omega$ is not simply connected, then $(f,\Omega) \in \mathbf{F}_0$ is the analytic continuation of a restriction to a simply connected subregion of $\Omega$, and since the restriction belongs to $\mathbf{F}$ so does $(f,\Omega)$ because of the property 1.

\index{Differential equation|)}\index{Function!analytic!global|)}\index{Function!hypergeometric|)}\index{Global analytic function|)}\index{Hypergeometric differential equation|)}\index{Linear differential equation|)}